\documentclass[11pt,a4paper]{article}

%% ── packages ──────────────────────────────────────────────────────
\usepackage[utf8]{inputenc}
\usepackage[T1]{fontenc}
\usepackage{amsmath,amssymb,amsthm,mathtools}
\usepackage{stmaryrd}            % \llbracket, \rrbracket
\usepackage{enumitem}
\usepackage{hyperref}
\usepackage[margin=1in]{geometry}
\usepackage{xcolor}
\usepackage{fancyhdr}
\usepackage{booktabs}
\usepackage{listings}

%% ── listings style ────────────────────────────────────────────────
\lstdefinestyle{rholang}{
  basicstyle=\small\ttfamily,
  keywordstyle=\bfseries,
  commentstyle=\itshape\color{gray},
  morekeywords={for,new,let,in,match},
  literate={<-}{{$\leftarrow$}}2
           {!}{{!}}1
           {|}{{$\mid$}}1,
  breaklines=true,
  frame=single,
  framesep=6pt,
  xleftmargin=1em,
  xrightmargin=1em,
  aboveskip=\medskipamount,
  belowskip=\medskipamount,
}

%% ── theorem environments ──────────────────────────────────────────
\newtheorem{theorem}{Theorem}[section]
\newtheorem{lemma}[theorem]{Lemma}
\newtheorem{proposition}[theorem]{Proposition}
\newtheorem{definition}[theorem]{Definition}
\newtheorem{remark}[theorem]{Remark}
\newtheorem{example}[theorem]{Example}

%% ── macros ────────────────────────────────────────────────────────
\newcommand{\lsem}{\llbracket}
\newcommand{\rsem}{\rrbracket}

% calculus syntax
\newcommand{\Par}{\mathbin{|}}
\newcommand{\forc}{\mathbf{for}}
\newcommand{\send}[2]{#1 !\!\left( #2 \right)}
\newcommand{\recv}[3]{\forc\!\left( #1 \leftarrow #2 \right) #3}
\newcommand{\nil}{\mathbf{0}}
\newcommand{\drop}[1]{{\ast} #1}
\newcommand{\quot}[1]{@#1}
\newcommand{\subst}[3]{#1\{#2 / #3\}}
\newcommand{\signed}[2]{\{#1\}_{#2}}

% signatures and tokens
\newcommand{\comp}{\mathbin{*}}
\newcommand{\hashf}{\texttt{\#}}

% structural equivalence and reduction
\newcommand{\Seq}{\equiv_{\!S}}
\newcommand{\Neq}{\equiv_{\!N}}
\newcommand{\red}{\rightsquigarrow}

% free names
\newcommand{\FN}{\mathrm{FN}}

%% ── headers ───────────────────────────────────────────────────────
\pagestyle{fancy}
\fancyhf{}
\fancyhead[L]{\small\textsc{F1R3FLY.io}}
\fancyhead[R]{\small\textsc{A Spectral Decomposition of Phlogiston}}
\fancyfoot[C]{\thepage}
\renewcommand{\headrulewidth}{0.4pt}

%% ── title ─────────────────────────────────────────────────────────
\title{%
  \textbf{Cost-Accounted Rho Calculus}\\[6pt]
  \large A Spectral Decomposition of Phlogiston}
\author{%
  L.\ Gregory Meredith\thanks{F1R3FLY.io; F1R3FLY Industries.
  \texttt{lgreg.meredith@f1r3fly.io}}%
}
\date{May 2026}

%% ══════════════════════════════════════════════════════════════════
\begin{document}
\maketitle

\begin{abstract}
We present a revision of the rho calculus in which cost
accounting is not an external extension to be translated back
into the core language, but is \emph{folded directly into the
grammar and operational semantics} of the calculus itself.
Signatures become a species of name alongside quoted processes;
token stacks become first-class terms; and the
for-comprehension and send operators range over \emph{signed
terms} rather than bare processes.  The result is a calculus in
which every communicating process must engage the cost layer
--- embodying the principle of \emph{valuable friction} at the
syntactic level.  The signature grammar is parameterized by a
type variable~$G$ representing the cryptographic backend,
making the calculus generic over signature schemes.

We identify the for-comprehension as the transaction primitive
of the calculus: the binding specification is the transaction,
the substitution is the witness, and application to the
continuation is the state update.  Individual transactions are
atomic by construction.  Multi-step transactions are made atomic
by a decidable static analysis at deployment-acceptance time: a
linear resource proof verifying that every for-comprehension in
the deployment's call graph is funded before any step executes.

Because the rho calculus is the formal core of Rholang, this
revised calculus serves as the design specification for
rearchitecting phlogiston accounting in Rholang~1.4 (MeTTaIL)
and the F1R3FLY platform.
\end{abstract}

\tableofcontents
\bigskip

%% ══════════════════════════════════════════════════════════════════
\section{Introduction}
\label{sec:intro}

The rho calculus~\cite{MeredithRadestock2005} is a reflective
higher-order process calculus: channels are the quoted codes of
processes, and processes may be recovered from channels by
dequotation.  As the formal core of
Rholang~\cite{RholangSpec}, the rho calculus underpins the
computation and storage model of the F1R3FLY platform.

Rholang's cost-accounting mechanism --- the \emph{phlogiston}
(phlo) system --- gates every deploy behind a token balance
associated with a digital signature.  In previous work, this
mechanism was specified as an \emph{extension} of the rho
calculus: additional syntax for signatures and tokens, and
additional rewrite rules for gated communication.  One could
then \emph{translate} the extension back into the pure rho
calculus, demonstrating that cost accounting is expressible
within the language.

This paper takes the next step.  Rather than treating cost
accounting as an external layer and then recovering it via
translation, we \emph{fold it directly into the calculus}.  The
key design moves are:

\begin{enumerate}[label=(\roman*)]
  \item \textbf{Signatures are names.}  The name sort is
    extended so that $x, y \;::=\; \quot{T} \mid s$, where~$s$
    ranges over signatures.  Signatures and quoted processes
    coexist as first-class channel names.
  \item \textbf{Two axes of reflection.}  Alongside the
    structural quoting operator~$\quot{T}$ (the code of a
    signed term used as a channel), a \emph{cryptographic
    quoting}
    operator~$\hashf P$ (the hash of a process) appears in the
    signature grammar.  This gives the calculus two axes of
    reflection: one for communication, one for authentication.
  \item \textbf{Signed terms pervade the syntax.}  The
    for-comprehension and send take \emph{signed terms}~$T, U$
    as their body and payload respectively, rather than bare
    processes.  One cannot write a communicating process without
    engaging the cost layer.
  \item \textbf{Token stacks are first-class terms.}  Token
    stacks $S \;::=\; () \mid s : S$ are included in the
    signed-term grammar, so they participate in reduction on
    equal footing with signed processes.
  \item \textbf{Parametric polymorphism.}  Parallel composition
    is overloaded: $P \Par Q$ composes processes, $T \Par U$
    composes signed terms.  The operator is parametrically
    polymorphic in the sort it composes.
  \item \textbf{Generic cryptography.}  The signature grammar
    is parameterized by a type variable~$G$ representing the
    ground signature scheme (Ed25519, secp256k1, etc.),  making
    the calculus independent of any particular cryptographic
    backend.
\end{enumerate}

\noindent
Because the rho calculus \emph{is} the core of Rholang, this
revised calculus is not a theoretical curiosity --- it is the
design specification for rearchitecting phlogiston accounting in
Rholang~1.4 (MeTTaIL) and the F1R3FLY platform.

%% ══════════════════════════════════════════════════════════════════
\section{Motivation: From Run-to-Completion to Concurrent Acceptance}
\label{sec:motivation}

The redesign presented in this paper is motivated by a
fundamental architectural limitation in the current F1R3FLY
implementation, inherited from RChain.

\subsection{The Run-to-Completion Model}
\label{sec:run-to-completion}

In the current architecture, each deployment submitted to a
F1R3Node is executed under \emph{run-to-completion} semantics:
the RSpace tuplespace engine accepts a deployment, executes it
to termination (or until phlogiston is exhausted), commits the
resulting state changes, and only then accepts the next
deployment.  This model introduces two related problems:

\begin{enumerate}[label=(\roman*)]
  \item \textbf{Implicit transaction boundaries.}  The entire
    deployment is treated as a single, implicit transaction.
    There is no formal specification of what constitutes a
    transaction --- the deployment boundary \emph{is} the
    transaction boundary, by convention rather than by
    design.  This conflates the unit of submission (the
    deployment) with the unit of atomicity (the transaction),
    making it impossible to express deployments that contain
    multiple independent transactions or to share
    transactional structure across deployments.

  \item \textbf{Limited availability.}  Because RSpace
    executes deployments sequentially, it is \emph{locked}
    for the duration of each deployment's execution.  No
    other deployment can be accepted, pattern-matched, or
    reduced while the current deployment is running.  The
    throughput of the F1R3Node is bounded by the execution
    time of the slowest deployment in the queue.  This is a
    severe bottleneck: a single expensive deployment blocks
    all other activity on the node.
\end{enumerate}

\subsection{The Merge Analysis Burden}
\label{sec:merge-burden}

To recover concurrency in the current architecture, F1R3FLY
(following RChain) employs \emph{merge analysis}: after
multiple deployments have been executed independently (e.g.,
by different validators or in speculative parallel
execution), their resulting states must be \emph{merged}
into a single consistent state.  The merge analysis
determines whether the independently computed states are
compatible (no conflicting writes to the same tuplespace
entries) and, if so, how to combine them.

Merge analysis is expensive, complex, and a persistent
source of bugs.  It requires tracking read and write sets
for every deployment, computing conflict graphs, and
resolving or rejecting conflicting state transitions.
Moreover, the analysis is \emph{post-hoc}: deployments are
executed speculatively, and conflicts are detected only
after the work has been done.  Wasted computation is the
norm, not the exception.

\subsection{The New Model: Acceptance by Linear Proof}
\label{sec:new-model}

The cost-accounted rho calculus eliminates both problems by
replacing run-to-completion with \emph{concurrent acceptance
gated by linear resource proofs}.

Under the new model:
\begin{enumerate}[label=(\roman*)]
  \item \textbf{Explicit transaction boundaries.}  The
    for-comprehension is the transaction primitive
    (Section~\ref{sec:for-as-txn}).  Transaction boundaries
    are determined by the syntactic structure of signed terms,
    not by deployment boundaries.  A deployment may contain
    multiple transactions; multiple deployments may
    participate in a single financial transaction (via shared
    signature channels).

  \item \textbf{Concurrent acceptance.}  RSpace can accept
    deployments \emph{at any time}, provided each deployment
    carries a linear resource proof
    (Section~\ref{sec:linear-proof}) demonstrating that
    sufficient tokens exist to complete all of its
    transactions.  Multiple deployments can be active
    simultaneously, each drawing from its own committed token
    supply.  RSpace is never locked by a single deployment.

  \item \textbf{Throughput.}  Because acceptance is gated by a
    static check (linear in the AST size) rather than by
    sequential execution, the F1R3Node can accept and begin
    executing deployments in parallel.  The throughput
    bottleneck shifts from ``one deployment at a time'' to
    ``as many deployments as there are non-conflicting token
    supplies'' --- a dramatic improvement.

  \item \textbf{Merge via signatures and cost accounting.}
    The need for post-hoc merge analysis is substantially
    reduced.  In the cost-accounted calculus, conflicts are
    detected \emph{at acceptance time} through the linear
    resource proof: two deployments that compete for the same
    tokens are two proof obligations competing for the same
    linear hypotheses, and at most one can succeed
    (Section~\ref{sec:deploy-boundaries}).  Merging is no
    longer a separate analysis pass over execution traces ---
    it is a consequence of the signature and cost-accounting
    structure of the calculus itself.

    Specifically, deployments signed by \emph{different}
    signatures draw from \emph{disjoint} token pools and
    cannot conflict.  They may be executed in parallel and
    their results composed without merge analysis.
    Deployments signed by the \emph{same} signature (or
    overlapping compound signatures) compete for the same
    token pool and are serialized by the linear proof ---
    again without merge analysis.  The only case requiring
    attention is deployments that interact via shared
    data channels, and even here the cost-accounting
    structure provides a natural serialization order.
\end{enumerate}

\begin{remark}[From blocking to budgeting]
  The architectural shift can be summarized in a single
  sentence: the current model asks ``is the tuplespace
  free?''\ (blocking); the new model asks ``is this
  deployment funded?''\ (budgeting).  The first question
  serializes all activity.  The second question can be
  answered independently for each deployment, enabling
  concurrency proportional to the diversity of the
  token supply.
\end{remark}

\subsection{New Capability: Multi-Signature Execution}
\label{sec:multisig-motivation}

Beyond streamlining execution and increasing throughput, the
cost-accounted calculus introduces a feature absent from the
current architecture: \emph{$n$-party multi-signature
execution}.

The compound signature $s_1 \comp s_2 \comp \cdots \comp s_n$
in the signed term
$\signed{P}{s_1 \comp \cdots \comp s_n}$
requires that \emph{all} $n$ parties explicitly sign off on the
execution of~$P$.  The rewrite rules enforce this structurally:
the process cannot fire unless tokens from every component
signature are present.  There is no separate multi-sig
verification layer, no threshold-signature scheme, no
off-chain coordination protocol.  Multi-sig is a consequence
of the grammar.

\paragraph{Programmable tokenized capabilities.}
The interaction between compound signatures and token stacks
yields a rich \emph{object-capability} model with
programmable amplification and attenuation:

\begin{description}[style=nextline, leftmargin=2em, nosep]
  \item[Attenuation.]
    A token stack limits what a capability can do.  A process
    signed by~$s$ with a two-token stack $s : s : ()$ can
    perform at most two funded actions.  By controlling the
    depth of the token stack granted to a subprocess, a parent
    process \emph{attenuates} the subprocess's authority.
    This is token-metered capability attenuation: the child
    can do no more than the parent has budgeted.

  \item[Amplification.]
    Composing signatures via $\comp$ \emph{amplifies}
    authority.  A process signed by $s_1 \comp s_2$ can
    perform actions that neither~$s_1$ nor~$s_2$ could
    authorize alone.  For example, a funds-transfer contract
    might require $\signed{\mathit{transfer}}
    {s_{\mathit{sender}} \comp s_{\mathit{compliance}}}$,
    meaning the transfer executes only if both the sender and
    the compliance officer have signed.  Neither party alone
    can move funds; their composed authority can.

  \item[Programmable delegation.]
    Because signatures are names and token stacks are
    first-class terms, a process can \emph{delegate}
    authority by sending tokens on a signature channel.  A
    manager process holding $s_{\mathit{mgr}} : S$ can send a
    portion of its stack to a subordinate process, granting
    the subordinate a metered budget of the manager's
    authority.  When the subordinate exhausts its budget, it
    halts --- no revocation protocol needed.

  \item[Just-in-time funding.]
    Because token stacks are first-class terms, they can be
    \emph{constructed and supplied dynamically} --- not merely
    pre-allocated at deployment time.  A contract can mint a
    token stack $s : s : ()$ and send it on a signature
    channel at the precise moment it is needed, funding a
    capability on an as-needed basis.

    This enables natural \emph{on-boarding} patterns: a
    platform contract grants a new user a small slush fund
    (a shallow token stack) to explore the system.  The
    user's initial deployments are funded by this grant,
    letting them experiment without acquiring tokens of
    their own.  Once the slush fund is exhausted, further
    execution requires the user to acquire tokens through
    the token economy --- purchase, earn, or receive via
    delegation.  The transition from subsidized exploration
    to self-funded participation is not a policy check
    external to the calculus; it is a \emph{resource
    exhaustion event} within it.
\end{description}

\noindent
This capability model is \emph{intrinsic} to the calculus:
amplification, attenuation, delegation, and just-in-time
funding are not library features or design patterns but
direct consequences of the signature algebra and token-stack
semantics.  The F1R3FLY platform's token-attenuated
object-capability model for agentic AI
governance~\cite{VACRADA} is a direct application of this
structure.

%% ══════════════════════════════════════════════════════════════════
\section{The Cost-Accounted Rho Calculus}
\label{sec:calculus}

\subsection{Syntax}
\label{sec:syntax}

The calculus has four mutually defined syntactic categories:
\emph{processes}, \emph{names}, \emph{signed terms}, and
\emph{signatures}.

\begin{definition}[Processes and names]
  \label{def:proc-name}
  \begin{align}
    P, Q &\;::=\; \nil
          \;\mid\; \recv{y}{x}{T}
          \;\mid\; \send{x}{U}
          \;\mid\; P \Par Q
          \;\mid\; \drop{x}
    \label{eq:proc-syntax}
    \\[4pt]
    x, y &\;::=\; \quot{T} \;\mid\; s
    \label{eq:name-syntax}
  \end{align}
\end{definition}

\noindent
The intended readings carry over from the pure rho calculus,
with one critical change: the for-comprehension
$\recv{y}{x}{T}$ has a \emph{signed term}~$T$ as its
continuation, and the send $\send{x}{U}$ carries a signed
term~$U$ as its payload.

\begin{definition}[Signed terms and token stacks]
  \label{def:signed-term}
  \begin{align}
    T, U &\;::=\; \signed{P}{s}
          \;\mid\; T \Par U
          \;\mid\; S
    \label{eq:signed-term-syntax}
    \\[4pt]
    S &\;::=\; ()
       \;\mid\; s : S
    \label{eq:token-syntax}
  \end{align}
\end{definition}

\noindent
A signed term is one of:
\begin{description}[style=nextline, leftmargin=2em]
  \item[$\signed{P}{s}$]
    A process~$P$ signed by signature~$s$.  The braces
    delimit the scope of the signature.
  \item[$T \Par U$]
    Parallel composition of signed terms --- the same
    $\Par$ operator as at the process level, lifted to
    signed terms.
  \item[$S$]
    A token stack, representing a phlogiston balance.
\end{description}

\noindent
A token stack $s_1 : s_2 : \cdots : s_n : ()$ is read as a
sequence of signature-indexed balance layers.  The empty
stack~$()$ represents a depleted balance.

\begin{definition}[Signatures, parameterized by~$G$]
  \label{def:sig}
  \begin{equation}
    s(G) \;::=\; g
          \;\mid\; \hashf P
          \;\mid\; s \comp s
    \label{eq:sig-syntax}
  \end{equation}
  where $g \in G$ is a \emph{ground signature} drawn from the
  cryptographic backend~$G$.
\end{definition}

\noindent
The three forms of signature are:
\begin{description}[style=nextline, leftmargin=2em]
  \item[$g$]
    A \emph{ground signature}: an opaque value from the
    cryptographic backend~$G$ (e.g., an Ed25519 public key, a
    secp256k1 key hash).  The calculus is parametrically
    polymorphic in~$G$.
  \item[$\hashf P$]
    \emph{Cryptographic quoting}: the hash of a process~$P$,
    yielding a signature.  This is the authentication-axis
    analog of structural quoting~$\quot{P}$.
  \item[$s_1 \comp s_2$]
    \emph{Signature composition}: the compound authority of
    $s_1$ and~$s_2$ acting jointly.
\end{description}

\begin{remark}[Two axes of reflection]
  \label{rem:two-axes}
  The calculus now has two quoting operators:
  \begin{center}
    \begin{tabular}{lll}
      \textbf{Operator} & \textbf{Axis} & \textbf{Purpose} \\
      \midrule
      $\quot{T}$ & Structural  & Signed term $\to$ channel name \\
      $\hashf P$ & Cryptographic & Process $\to$ signature \\
    \end{tabular}
  \end{center}
  Both are forms of \emph{reflection}: they take a computational
  entity and produce a \emph{name} (a datum that can be
  communicated, compared, and used as an address).  The
  difference is that $\quot{T}$ preserves full structural
  information --- including the signature and cost provenance
  of~$T$, and the signed term can be recovered
  via~$\drop{x}$ --- while $\hashf P$ is a one-way,
  collision-resistant digest, suitable for authentication but
  not for recovery.
\end{remark}

\subsection{Parametric Polymorphism of Parallel Composition}
\label{sec:par-poly}

The parallel composition operator~$\Par$ appears at two levels:

\begin{enumerate}[nosep]
  \item $P \Par Q$ --- parallel composition of processes.
  \item $T \Par U$ --- parallel composition of signed terms.
\end{enumerate}

\noindent
This is a parametrically polymorphic operator: the same symbol,
the same algebraic structure (commutative monoid with
appropriate identity), but ranging over different sorts.  At the
process level the identity is~$\nil$; at the signed-term level
the identity is the empty stack~$()$.

\subsection{Free Names}
\label{sec:fn}

The binding discipline carries over from the pure rho calculus:
$\recv{y}{x}{T}$ binds~$y$.  Free names are computed
recursively over both processes and signed terms:

\begin{align*}
  \FN(\nil) &= \emptyset \\
  \FN(\recv{y}{x}{T})
    &= \bigl(\FN(T) \setminus \{y\}\bigr) \cup \{x\} \\
  \FN(\send{x}{U})
    &= \FN(U) \cup \{x\} \\
  \FN(P \Par Q) &= \FN(P) \cup \FN(Q) \\
  \FN(\drop{x}) &= \{x\}
\end{align*}

\noindent
Extended to signed terms:
\begin{align*}
  \FN(\signed{P}{s})
    &= \FN(P) \cup \FN_s(s) \\
  \FN(T \Par U)
    &= \FN(T) \cup \FN(U) \\
  \FN(())   &= \emptyset \\
  \FN(s : S)
    &= \FN_s(s) \cup \FN(S)
\end{align*}

\noindent
where $\FN_s$ extracts names from signatures:
\begin{align*}
  \FN_s(g) &= \emptyset \\
  \FN_s(\hashf P) &= \FN(P) \\
  \FN_s(s_1 \comp s_2)
    &= \FN_s(s_1) \cup \FN_s(s_2)
\end{align*}

\subsection{Structural Equivalence}
\label{sec:struct-equiv}

Structural equivalence~$\Seq$ is the smallest equivalence
relation on processes including $\alpha$-equivalence and making
$(P, \Par, \nil)$ a commutative monoid:

\begin{align*}
  P \Par \nil    &\Seq P           \\
  P \Par Q       &\Seq Q \Par P    \\
  (P \Par Q) \Par R &\Seq P \Par (Q \Par R)
\end{align*}

\noindent
Likewise, at the signed-term level, $(T, \Par, ())$ forms a
commutative monoid under~$\Seq$:

\begin{align*}
  T \Par ()      &\Seq T           \\
  T \Par U       &\Seq U \Par T    \\
  (T \Par U) \Par V &\Seq T \Par (U \Par V)
\end{align*}

\subsection{Substitution}
\label{sec:subst}

Substitution $T\{\quot{U}/y\}$ is the standard capture-avoiding
substitution of~$\quot{U}$ for~$y$ in~$T$, extended to signed
terms.  Following the pure rho calculus, the only non-trivial
case is dequotation:
\[
  (\drop{x})\{\quot{U}/y\}
  \;=\;
  \begin{cases}
    U & \text{if } y \Neq x, \\
    \drop{x} & \text{otherwise,}
  \end{cases}
\]
where $\Neq$ denotes name equality.

\begin{remark}[Signed-term substitution]
  \label{rem:signed-subst}
  Because the name sort is $x, y \;::=\; \quot{T} \mid s$,
  the substitution $T\{\quot{U}/y\}$ naturally substitutes the
  quoted \emph{signed term}~$U$, not a bare process extracted
  from~$U$.  There is no unwrapping.  This means the signature
  and cost information associated with~$U$ is \emph{preserved
  through communication}: a signed term received on a channel
  retains its provenance.
\end{remark}

\subsection{Rewrite Rules}
\label{sec:rewrites}

The operational semantics consists of five rewrite rules, each
a variation on the \textsc{Comm} rule of the pure rho calculus
gated by token consumption.  In every rule, a communication
fires only when matching tokens are present, and the
consumed tokens are released as residual terms.

\bigskip
\noindent
\textbf{Rule 1} --- Single signature, whole redex:
\begin{equation}
  \boxed{
    \signed{\recv{y}{x}{T} \Par \send{x}{U}}{s}
    \;\Par\; s : S
    \;\;\red\;\;
    \subst{T}{\quot{U}}{y} \Par S
  }
  \label{eq:rule1}
\end{equation}

\noindent
The entire redex (input and output) is signed by~$s$.  A single
layer of token is consumed; the stack remainder~$S$ is released.

\bigskip
\noindent
\textbf{Rule 2} --- Compound signature, whole redex, split
tokens:
\begin{equation}
  \boxed{
    \signed{\recv{y}{x}{T} \Par \send{x}{U}}{s_1 \comp s_2}
    \;\Par\; s_1 : S \Par s_2 : S'
    \;\;\red\;\;
    \subst{T}{\quot{U}}{y} \Par S \Par S'
  }
  \label{eq:rule2}
\end{equation}

\noindent
The redex is signed by the compound authority~$s_1 \comp s_2$.
Both component signatures must supply tokens; both remainders
are released.

\bigskip
\noindent
\textbf{Rule 3} --- Compound signature, whole redex, combined
token:
\begin{equation}
  \boxed{
    \signed{\recv{y}{x}{T} \Par \send{x}{U}}{s_1 \comp s_2}
    \;\Par\; (s_1 \comp s_2) : S
    \;\;\red\;\;
    \subst{T}{\quot{U}}{y} \Par S
  }
  \label{eq:rule3}
\end{equation}

\noindent
The compound authority~$s_1 \comp s_2$ holds a single combined
token.  One layer is consumed; the remainder~$S$ is released.

\bigskip
\noindent
\textbf{Rule 4} --- Split processes, split tokens:
\begin{equation}
  \boxed{
    \signed{\recv{y}{x}{T}}{s_1}
    \;\Par\; \signed{\send{x}{U}}{s_2}
    \;\Par\; s_1 : S \Par s_2 : S'
    \;\;\red\;\;
    \subst{T}{\quot{U}}{y} \Par S \Par S'
  }
  \label{eq:rule4}
\end{equation}

\noindent
Input and output are signed separately.  Each signer supplies
their own token.

\bigskip
\noindent
\textbf{Rule 5} --- Split processes, combined token:
\begin{equation}
  \boxed{
    \signed{\recv{y}{x}{T}}{s_1}
    \;\Par\; \signed{\send{x}{U}}{s_2}
    \;\Par\; (s_1 \comp s_2) : S
    \;\;\red\;\;
    \subst{T}{\quot{U}}{y} \Par S
  }
  \label{eq:rule5}
\end{equation}

\noindent
Input and output are signed separately, but the token is held
jointly under the compound authority.

\bigskip
\noindent
As in the pure rho calculus, these rules are closed under
parallel context and structural equivalence:
\begin{gather}
  \frac{P \red P'}{P \Par Q \red P' \Par Q}
  \qquad\text{(\textsc{Par})}
  \label{eq:par-ctx}
  \\[6pt]
  \frac{P \Seq P' \qquad P' \red Q' \qquad Q' \Seq Q}
       {P \red Q}
  \qquad\text{(\textsc{Struct})}
  \label{eq:struct-ctx}
\end{gather}

\subsection{Summary of the Five Rules}
\label{sec:rule-summary}

The five rules can be organized along two axes:

\begin{center}
  \begin{tabular}{l|cc}
    & \textbf{Split tokens} & \textbf{Combined token} \\
    \midrule
    \textbf{Whole redex, single sig}
      & Rule~1 & Rule~1 \\
    \textbf{Whole redex, compound sig}
      & Rule~2 & Rule~3 \\
    \textbf{Split processes}
      & Rule~4 & Rule~5 \\
  \end{tabular}
\end{center}

\noindent
Every rule follows the same pattern: \emph{communication is
gated by token consumption}.  The variations arise from the
granularity of signing (whole redex vs.\ split processes) and
the granularity of token holding (split vs.\ combined).

\subsection{Syntactic Sugar}
\label{sec:sugar}

Two abbreviations streamline the common case of signing a
for-comprehension together with its continuation.

\begin{definition}[Uniform signing]
  \label{def:sugar-uniform}
  \begin{equation}
    \signed{\recv{y}{x}{P}}{s}
    \;\;\triangleq\;\;
    \signed{
      \recv{y}{x}{\signed{P}{s}}
    }{s}
    \label{eq:sugar-uniform}
  \end{equation}
\end{definition}

\noindent
When a for-comprehension and its continuation are signed by
the \emph{same} signature~$s$, the abbreviation on the left
expands to two nested signed layers on the right.  The outer
$\signed{\cdot}{s}$ pays for the transaction (rendezvous and
matching); the inner $\signed{\cdot}{s}$ pays for executing
the continuation.  This is the common case for single-party
execution: both the communication and its effect are
authorized and funded by the same signer.

\begin{definition}[Linear transfer of authority ($\multimap$)]
  \label{def:sugar-lollipop}
  \begin{equation}
    \signed{\recv{y}{x}{P}}{s_1 \multimap s_2}
    \;\;\triangleq\;\;
    \signed{
      \recv{y}{x}{\signed{P}{s_2}}
    }{s_1}
    \label{eq:sugar-lollipop}
  \end{equation}
\end{definition}

\noindent
The \emph{lollipop} notation
$\signed{\cdot}{s_1 \multimap s_2}$ denotes a
\emph{transfer of authority}: signature~$s_1$ authorizes and
funds the communication, but signature~$s_2$ authorizes and
funds the continuation.  The $\multimap$ is deliberately
chosen to evoke the linear implication of linear logic:
$s_1$'s resource is consumed to establish the communication,
and $s_2$'s resource funds the resulting computation.

\begin{example}[Cross-party service invocation]
  \label{ex:cross-party}
  A client~$c$ invokes a service operated by provider~$p$.
  The client pays for the call; the provider pays for the
  execution:
  \[
    \signed{
      \recv{\mathit{result}}{\mathit{service}}{\mathit{handleResult}}
    }{c \multimap p}
    \;\;=\;\;
    \signed{
      \recv{\mathit{result}}{\mathit{service}}{
        \signed{\mathit{handleResult}}{p}
      }
    }{c}
  \]
  The client's token funds the rendezvous with the service
  channel; the provider's token funds the execution of
  $\mathit{handleResult}$.  Neither party can impose costs
  on the other beyond what the sugar makes explicit.
\end{example}

\begin{remark}[Composability]
  Both sugar forms compose with compound signatures.  For
  example,
  $\signed{\recv{y}{x}{P}}{(s_1 \comp s_2) \multimap s_3}$
  requires both~$s_1$ and~$s_2$ to authorize the
  communication, while $s_3$ alone funds the continuation.
  Chains of transfers are also expressible:
  nested for-comprehensions can each carry their own
  $\multimap$, threading authority through a pipeline of
  parties.
\end{remark}

%% ══════════════════════════════════════════════════════════════════
\section{Design Analysis}
\label{sec:analysis}

\subsection{Signatures as Names: Eliminating the Translation Layer}
\label{sec:sig-as-name}

In a previous formulation, the cost-accounted calculus was
specified as an extension of the pure rho calculus, and a
compositional translation $\mathcal{N}\lsem\cdot\rsem$ mapped
signatures to names, tokens to messages, and signed processes
to fuel-gated wrappers.  While that translation demonstrated
expressibility, it introduced an indirection: the
implementation had to maintain two representations and a mapping
between them.

The present formulation eliminates this indirection.  By
declaring $x, y \;::=\; \quot{T} \mid s$, signatures \emph{are}
names.  There is nothing to translate.  The implementation
maintains a single namespace in which quoted signed terms and
signatures coexist.  This has direct consequences for the
Rholang runtime:

\begin{itemize}[nosep]
  \item The tuplespace (RSpace) stores both data channels
    ($\quot{T}$) and signature channels ($s$) in the same
    structure.
  \item Pattern matching in for-comprehensions can match on
    both kinds of name, enabling contracts that dispatch on
    the kind of authority presented.
  \item Namespace isolation is enforced by the existing
    namespace logic~\cite{MeredithRadestock2005}, now extended
    to cover the signature sort.
\end{itemize}

\subsection{Cryptographic Quoting}
\label{sec:crypto-quote}

The operator $\hashf P$ introduces a second axis of reflection.
If structural quoting~$\quot{T}$ is the ``mirror'' that
reflects a signed term into a name preserving all structure
(including signature and cost provenance), then
cryptographic quoting~$\hashf P$ is a ``one-way mirror'': it
produces a name (usable as a signature, an address, a key) from
which the original process cannot be recovered.

This distinction is fundamental.  A structurally quoted signed
term can be dequoted ($\drop{(\quot{T})} = T$) and the process
within it re-executed.  A
cryptographically quoted process cannot: $\hashf P$ is a
commitment, not a capability.  This makes $\hashf P$ suitable
for authentication, content addressing, and integrity
verification --- all operations that require binding to content
without granting the power to execute it.

In the F1R3FLY platform, $\hashf P$ will be implemented as a
configurable hash function (SHA-256, Blake2b, etc.) applied to
the serialized process tree.  The parameterization of~$s$ by~$G$
ensures that the choice of hash function is a deployment
decision, not a calculus-level commitment.

\subsection{Valuable Friction at the Syntactic Level}
\label{sec:valuable-friction}

The F1R3FLY platform's design philosophy of \emph{valuable
friction} holds that every resource-consuming operation should
carry an explicit cost, and that this cost is not merely a
runtime check but a \emph{structural feature} of the language.
The revised calculus embodies this principle:

\begin{itemize}[nosep]
  \item A for-comprehension $\recv{y}{x}{T}$ cannot fire unless
    its continuation~$T$ includes (possibly nested) signed
    terms with matching tokens.
  \item A send $\send{x}{U}$ carries a signed term~$U$ as
    payload, so the receiver inherits the sender's cost
    commitments.
  \item There is no way to write a communicating process that
    ``escapes'' the cost layer, because the grammar does not
    allow bare processes in the positions where communication
    occurs.
\end{itemize}

\subsection{Token Stacks as First-Class Terms}
\label{sec:token-stacks}

Token stacks $S \;::=\; () \mid s : S$ are lists of signatures,
representing layered balances.  Their inclusion in the
signed-term grammar ($T \;::=\; \cdots \mid S$) means that
tokens are not inert metadata --- they are terms that appear in
parallel compositions, flow through channels as parts of signed
messages, and participate in reduction.

This design enables several patterns:

\begin{itemize}[nosep]
  \item \textbf{Token delegation:}  A contract can receive a
    token stack on a channel and use it to fund subsequent
    communications.
  \item \textbf{Token splitting:}  A stack $s_1 : s_2 : S$ can
    be destructured by separate rules consuming $s_1$ and $s_2$
    independently.
  \item \textbf{Token markets:}  Because tokens are terms, they
    can be exchanged, auctioned, and composed using standard
    Rholang contract patterns.
\end{itemize}

\subsection{Genericity over the Cryptographic Backend}
\label{sec:genericity}

The signature grammar is parameterized: $s(G) \;::=\; g \mid
\hashf P \mid s \comp s$, where $g \in G$ and $G$ is a type
parameter.  This makes the entire calculus --- and therefore
the Rholang implementation --- generic over the cryptographic
backend.

Instantiating~$G$ yields a concrete calculus:
\begin{itemize}[nosep]
  \item $G = \mathrm{Ed25519}$:  ground signatures are Ed25519
    public keys.
  \item $G = \mathrm{secp256k1}$:  ground signatures are
    Ethereum-compatible key hashes.
  \item $G = \mathrm{Ed25519} + \mathrm{secp256k1}$:  a
    coproduct backend supporting both schemes.
\end{itemize}

\noindent
The $\comp$ operator composes signatures regardless of their
ground type, enabling cross-scheme multi-sig patterns.  The
implementation must provide a ground-signature equality
predicate and a hash function for each instantiation of~$G$;
the calculus itself is agnostic.

\subsection{From Phlogiston to Spectrum: The Decomposition of
  Cost}
\label{sec:spectrum}

In the original Rholang architecture, \emph{phlogiston} was a
singular, undifferentiated resource: every deploy consumed units
of the same universal fuel, regardless of who signed it or what
it computed.  The cost-accounted calculus decomposes this
singular resource into a \emph{spectrum} of signature-indexed
resources --- as many distinct fuels as there are economically
empowered actors.

Two metaphors illuminate this decomposition:

\paragraph{The prism.}
Phlogiston is white light.  The signature algebra acts as a
prism, splitting the undifferentiated beam into its many
constituent frequencies.  Each signature~$s$ indexes a distinct
spectral line --- a resource pool visible only to processes
signed by~$s$.  Compound signatures $s_1 \comp s_2$ produce
interference patterns: new frequencies that emerge only when
two spectral components are superposed.  The linear resource
proof at deployment-acceptance time is a \emph{spectral
analysis}: it verifies that every frequency the deployment
requires is present with sufficient intensity.

\paragraph{Algorithmic chemistry.}
Alternatively, one may view the calculus as an \emph{algorithmic
chemistry} in which economic agency is the catalyst for
computation.  Processes are reagents; tokens are catalytic
substrates; and the rewrite rules are reaction rules that fire
only when the appropriate catalyst is present.  A signed
process $\signed{P}{s}$ is a reaction waiting for the
catalyst~$s$.  A token stack $s : S$ is a quantity of catalytic
substrate.  Communication is a catalyzed reaction: the substrate
is consumed, the reagents are transformed, and the products
(the continuation plus the remaining substrate) are released.

In this reading, the diversity of signatures corresponds to the
diversity of catalysts in a chemical system.  Different
reactions require different catalysts.  A compound
signature~$s_1 \comp s_2$ is a \emph{co-catalyst}: both
substances must be present for the reaction to proceed.  The
economy of the F1R3FLY platform is an algorithmic chemistry in
which the diversity and abundance of catalysts determine the
rate and variety of computation.

\begin{remark}[Phlogiston as a limit case]
  The original phlogiston model is recovered by assigning all
  actors the same signature~$s_0$.  In the prism metaphor, this
  collapses the spectrum back to white light.  In the chemistry
  metaphor, this reduces the system to a single universal
  catalyst.  The cost-accounted calculus generalizes from this
  monochromatic / mono-catalytic limit to a full spectrum of
  economic agency.
\end{remark}

\subsection{Funding Slots via Unforgeable Channels}
\label{sec:funding-slots}

The linear-transfer sugar
$\signed{\recv{y}{x}{P}}{s_1 \multimap s_2}$ requires the
continuation's funder~$s_2$ to be named at compile time.  In
practice, one often wants to write a contract whose
continuation can be funded by \emph{any} party willing to pay
--- without knowing their identity in advance.  Adding
existential quantification over signatures would achieve this
but at the cost of complicating the linear resource proofs
(the demand function~$\Delta_s$ would need to range over
unknown signatures).

A simpler and more natural solution exploits a feature that
Rholang already provides: \emph{unforgeable channels} via
$\mathbf{new}$.  Rather than existentially quantifying over
the funder's signature, a contract creates a fresh funding
channel --- a \emph{funding slot} --- and listens for tokens
on it.  Anyone who knows the channel can deposit tokens there.
The linear proof at acceptance time checks whether tokens are
present on the funding slot, not who deposited them.

\begin{definition}[Funding slot]
  \label{def:funding-slot}
  A funding slot is a pattern in which a contract creates an
  unforgeable channel and uses it as the signature for a
  continuation:
  \[
    \mathbf{new}\;\mathit{slot}\;\mathbf{in}\;
    \bigl(
      \signed{\recv{y}{x}{P}}
        {s_1 \multimap \mathit{slot}}
      \;\Par\;
      \send{\quot{\mathit{slot}}}{(\cdots)}
    \bigr)
  \]
  where the send on~$\quot{\mathit{slot}}$ publishes the
  funding-slot address to potential funders.
\end{definition}

\noindent
The funding slot decouples three concerns:

\begin{description}[style=nextline, leftmargin=2em, nosep]
  \item[Authorization (who may initiate).]
    Controlled by $s_1$: only the signer of~$s_1$ can
    trigger the communication.
  \item[Funding (who pays for the continuation).]
    Open: anyone who deposits tokens on~$\mathit{slot}$.
  \item[Execution (what happens).]
    Determined by~$P$: the continuation runs once the
    slot is funded.
\end{description}

\begin{example}[Sponsored execution]
  \label{ex:sponsored}
  A new user~$u$ wants to run a deployment~$D$ but holds no
  tokens.  A sponsor~$\sigma$ is willing to fund the
  execution.  The on-boarding contract creates a funding slot:
  \begin{align*}
    \mathbf{new}\;\mathit{slot}\;\mathbf{in}\;
    \bigl(\,
      &\signed{\recv{\mathit{dep}}{u}{D}}
        {u \multimap \mathit{slot}}
      \\[-2pt]
      &\;\Par\;
      \send{\quot{\mathit{slot}}}
        {\mathit{slot} : \mathit{slot} : ()}
    \,\bigr)
  \end{align*}

  \noindent
  The sponsor funds the slot by depositing tokens:
  \[
    \send{\Sigma\lsem\mathit{slot}\rsem}{
      \mathit{slot} : \mathit{slot} : ()
    }
  \]
  The user~$u$ authorizes and triggers the communication
  (consuming $u$-tokens); the sponsor's tokens on the
  funding slot pay for the continuation.  The linear proof
  sees concrete tokens on a concrete channel --- no
  existential quantification is needed.

  When the slush fund is exhausted, the user must either
  attract another sponsor or acquire tokens of their own.
  The transition from subsidized to self-funded participation
  is a resource-exhaustion event, not a policy decision.
\end{example}

\begin{remark}[Unforgeability and security]
  The $\mathbf{new}$ operator guarantees that the funding
  slot is \emph{unforgeable}: only processes that have been
  explicitly given the channel name can deposit tokens on it.
  This prevents unauthorized parties from draining the slot
  or injecting spurious tokens.  The capability-security
  properties of unforgeable channels, already well-understood
  in the rho calculus, carry over directly to funding slots.
\end{remark}

%% ══════════════════════════════════════════════════════════════════
\section{Relationship to the Pure Rho Calculus}
\label{sec:relationship}

The cost-accounted calculus is a \emph{conservative extension}
of the pure rho calculus.  Setting $T = \signed{P}{s_0}$ for a
distinguished ``free'' signature~$s_0$ with unbounded token
supply, and collapsing the signed-term and process levels,
recovers the pure rho calculus exactly:

\[
  \signed{\recv{y}{x}{\signed{P'}{s_0}} \Par
    \send{x}{\signed{Q}{s_0}}}{s_0}
  \;\Par\; s_0 : S
  \;\;\red\;\;
  \signed{P'}{s_0}\{\quot{\signed{Q}{s_0}}/y\} \Par S
\]

\noindent
With $s_0 : S$ always available (supplied by a persistent
infrastructure process), this reduces to the standard
\textsc{Comm} rule.  The free signature~$s_0$ plays the role of
an ``always-funded'' system account, modeling the behavior of
the pure calculus where cost is not a concern.

This embedding also clarifies the implementation path: the
Rholang~1.4 runtime need only implement the cost-accounted
calculus, and the ``free'' mode (for testing, development, and
system-level contracts) is obtained by fixing~$s_0$ and
providing an inexhaustible token supply.

%% ══════════════════════════════════════════════════════════════════
\section{Implications for the F1R3FLY Platform}
\label{sec:implications}

\subsection{Rholang 1.4 (MeTTaIL) Integration}

MeTTaIL~\cite{MeTTaIL}, the Rholang~1.4 language, provides
means to define languages as triples of types (syntactic
categories), equations on types, and rewrites on types.  The
cost-accounted rho calculus is itself such a triple:

\begin{itemize}[nosep]
  \item \textbf{Types:}  $P$, $T$, $S$, $s(G)$, and the name
    sort.
  \item \textbf{Equations:}  Structural equivalence ($\Seq$)
    at both the process and signed-term levels.
  \item \textbf{Rewrites:}  Rules~1--5 and the contextual
    closure rules.
\end{itemize}

\noindent
This means the cost-accounted calculus can be \emph{defined
within MeTTaIL as a GSLT} (graph-structured lambda theory),
inheriting MeTTaIL's tooling for type checking, bisimulation,
and OSLF-generated logics.

\subsection{RSpace and the Tuplespace}

In the current F1R3FLY architecture, RSpace is a tuple-space
storage engine keyed by channel names.  The revised calculus
requires RSpace to store entries keyed by both~$\quot{T}$
and~$s$.  Since both are names, this is a matter of extending
the serialization format to accommodate the signature sort,
not a structural change to RSpace.

Token stacks~$S$ are stored as values in signature-keyed
entries.  The five rewrite rules correspond to tuplespace
match-and-consume operations that the reducer already performs,
extended with a signature-matching predicate.

\subsection{Casper CBC Consensus}

The cost-accounting rules interact with Casper CBC consensus at
the validator level:  a proposed block is valid only if every
communication it contains is backed by matching tokens.  The
revised calculus makes this validity check a syntactic
predicate: the block's contents must be well-typed in the
cost-accounted grammar, with token stacks present for every
signed communication.

\subsection{Implementation Path}

\begin{enumerate}[nosep]
  \item \textbf{Phase 1:}  Extend the Rholang parser and AST
    to include the signed-term, token-stack, and parameterized
    signature sorts.  Implement the five rewrite rules in the
    reducer.
  \item \textbf{Phase 2:}  Extend RSpace serialization to
    handle signature-keyed entries and token-stack values.
  \item \textbf{Phase 3:}  Define the cost-accounted calculus as
    a GSLT in MeTTaIL, enabling formal verification via OSLF.
  \item \textbf{Phase 4:}  Expose the signature and token APIs
    to user contracts, enabling composable metering,
    cross-chain fee markets, and delegated funding patterns.
\end{enumerate}

%% ══════════════════════════════════════════════════════════════════
\section{Transactions, Atomicity, and Static Analysis}
\label{sec:transactions}

The cost-accounted calculus reveals a precise correspondence
between the for-comprehension and the concept of a
\emph{transaction}.  This section develops that correspondence,
identifies a partial-funding vulnerability in multi-step
transactions, and shows that the vulnerability is eliminated by
a decidable static analysis at deployment-acceptance time.

\subsection{The For-Comprehension as Transaction Primitive}
\label{sec:for-as-txn}

Every for-comprehension $\recv{y}{x}{T}$ in the cost-accounted
calculus is a \emph{transaction} in the precise sense that it
decomposes into two funded computational actions:

\begin{enumerate}[nosep]
  \item \textbf{Rendezvous at the channel.}  The runtime must
    locate a matching send $\send{x}{U}$ on channel~$x$.  This
    is a search operation over the tuplespace, and its cost is
    paid by one token from the signer's stack.
  \item \textbf{Matching.}  The runtime must verify that the
    pattern~$y$ matches the offered payload~$U$ and compute
    the substitution $\{\quot{U}/y\}$.  This is a
    pattern-matching and substitution operation, and its cost
    is paid by a second token.
\end{enumerate}

\noindent
The \emph{transaction} is exactly what appears between the
round braces of the for-comprehension: the binding
specification $(y \leftarrow x)$.  Everything else --- the
continuation~$T$, the surrounding signed-term structure, the
token stack --- is the \emph{context} in which the transaction
executes.

\begin{definition}[Transaction, witness, state update]
  \label{def:txn}
  Given a signed for-comprehension
  $\signed{\recv{y}{x}{T}}{s}$ and a matching send
  $\send{x}{U}$:
  \begin{description}[style=nextline, leftmargin=2em, nosep]
    \item[Transaction]
      The binding specification $(y \leftarrow x)$: the
      declaration that we seek a value on channel~$x$ and
      will bind it to~$y$.
    \item[Witness]
      The substitution $\{\quot{U}/y\}$: the proof that the
      transaction succeeded.  A message was present on~$x$, it
      matched the pattern, and here is the concrete value that
      was bound.
    \item[State update]
      The application of the substitution to the
      continuation, $\subst{T}{\quot{U}}{y}$: the
      computational effect of the transaction on the system
      state.
    \item[Atomicity]
      The substitution is either applied or not.  There is no
      intermediate state.  The rewrite rule
      \[
        \signed{\recv{y}{x}{T} \Par \send{x}{U}}{s}
        \Par s : S
        \;\red\;
        \subst{T}{\quot{U}}{y} \Par S
      \]
      fires as a single, indivisible step.  If the token is
      absent or the send is missing, nothing happens ---
      the process blocks.
  \end{description}
\end{definition}

\begin{remark}[Database atomicity by construction]
  \label{rem:db-atomicity}
  It is important to distinguish two levels of atomicity that
  the community frequently conflates:
  \emph{database-transaction atomicity} and
  \emph{financial-transaction atomicity}.

  \textbf{Database-transaction atomicity} is the property of a
  single communication step: the substitution is either applied
  or not, the token is either consumed or not, the state is
  either updated or not.  In the cost-accounted rho calculus,
  this is a \emph{theorem} of the operational semantics ---
  the rewrite rules are defined so that a communication either
  fires completely or does not fire at all.  There is no rule
  that performs a ``partial'' communication.  No locks,
  two-phase commit, or rollback logs are needed.

  \textbf{Financial-transaction atomicity} is a different
  property: a multi-step operation (e.g., a payment decomposed
  into debit and credit) should either complete entirely or not
  begin at all.  This is \emph{not} guaranteed by the rewrite
  rules alone --- it requires the static analysis developed in
  Section~\ref{sec:linear-proof} and the deployment-boundary
  semantics of Section~\ref{sec:deploy-boundaries}.
\end{remark}

\subsection{The Partial Funding Vulnerability}
\label{sec:partial-funding}

Single-step transactions are atomic by construction.  The
vulnerability arises in \emph{multi-step} transactions: a
deployment that chains two or more for-comprehensions
sequentially may have enough tokens to fund the first step
but not the second.

Consider a payment from Alice to Bob, decomposed into two
operations:

\begin{enumerate}[nosep]
  \item \textbf{Debit:}  Extract a purse of amount~$n$ from
    Alice's wallet.
  \item \textbf{Credit:}  Deposit the purse into Bob's wallet.
\end{enumerate}

\noindent
Each operation is a for-comprehension, each wrapped in a
signed term $\signed{\cdots}{\mathit{Alice}}$, each consuming
tokens from Alice's stack.  If Alice's stack has enough tokens
for the debit but not the credit, the system reaches an
inconsistent intermediate state: funds extracted from Alice's
wallet, sitting in a purse, but the deposit into Bob's wallet
never fires because the fuel ran out.

This is a \emph{partial funding vulnerability}: the
transaction is semantically atomic (it should either complete
or not happen) but operationally non-atomic (it can stall
midway through).

\subsection{Rholang 1.2 Illustration}
\label{sec:rholang-illustration}

We illustrate the vulnerability using Rholang~1.2 syntax,
where the $\mathtt{?!}$ operator provides synchronous
(function-call-like) communication: $\mathtt{x?!(args)}$ sends
on~$x$ with a return channel and blocks until the reply.
All three code fragments below use the uniform-signing sugar
(Definition~\ref{def:sugar-uniform}):
$\signed{\recv{y}{x}{P}}{s}$ expands to
$\signed{\recv{y}{x}{\signed{P}{s}}}{s}$, so each signed
for-comprehension consumes \emph{two} tokens (one for the
communication, one for the continuation).

\medskip
\noindent
\textbf{Debit handler} --- extracts a purse from Alice's
wallet:
\begin{lstlisting}[style=rholang]
// Sugar: costs 2 Alice-tokens (comm + continuation)
{ for( amt <- aliceDebit?!( request, ... ) ){
     let ans = extractPurse( request, ... )
     in rtn!( ans )
} }_{ Alice }
\end{lstlisting}

\noindent
\textbf{Credit handler} --- deposits a purse into Bob's
account:
\begin{lstlisting}[style=rholang]
// Sugar: costs 2 Alice-tokens (comm + continuation)
{ for( amt <- bobCredit?!( ... ) ){
    new rcpt in
       addPurseAcct( amt, ... )
       in rtn!( rcpt )
} }_{ Alice }
\end{lstlisting}

\noindent
\textbf{Alice's token stack:}
\begin{lstlisting}[style=rholang]
Alice : Alice : Alice : Alice : Alice : Alice : S
\end{lstlisting}

\noindent
\textbf{Orchestrating deployment} --- chains the debit and
credit:
\begin{lstlisting}[style=rholang]
// Sugar: costs 2 Alice-tokens (comm + continuation)
{ for( amt <- aliceDebit?!( request, ... );
       rcpt <- bobCredit?!( amt, ... ) ){
     P
} }_{ Alice }
\end{lstlisting}

\noindent
The total token demand of this deployment is:
\begin{center}
  \begin{tabular}{lccc}
    \textbf{Signed layer}
      & \textbf{Comm}
      & \textbf{Cont.}
      & \textbf{Total} \\
    \midrule
    Orchestrator's $\signed{\cdots}{\mathit{Alice}}$
      & 1 & 1 & 2 \\
    Debit handler's $\signed{\cdots}{\mathit{Alice}}$
      & 1 & 1 & 2 \\
    Credit handler's $\signed{\cdots}{\mathit{Alice}}$
      & 1 & 1 & 2 \\
    \midrule
    \textbf{Total} & \textbf{3} & \textbf{3} & \textbf{6} \\
  \end{tabular}
\end{center}

\noindent
If Alice's stack has at least six tokens, all three
signed layers (and their continuations) are funded and the
transaction completes.  If the stack has only four tokens,
the orchestrator and the debit handler fire (consuming
four tokens), but the credit handler's signed layer finds
no token and \emph{blocks}.
Funds are stranded in the purse.

\subsection{Decomposing the Transaction Cost}
\label{sec:txn-cost-decomp}

The six-token count of Section~\ref{sec:rholang-illustration}
reflects the uniform-signing sugar expansion of three explicit
$\signed{\cdots}{\mathit{Alice}}$ layers.  However, the
$\mathtt{?!}$ operator itself desugars to a send followed by a
for-comprehension on the return channel, so each $\mathtt{?!}$
call involves a for-comprehension on \emph{each} side --- the
caller's and the handler's.  These internal for-comprehensions
are also under signed scope and therefore also consume tokens.

A full accounting, decomposed into the two computational
actions of Section~\ref{sec:for-as-txn}, is:

\begin{center}
  \begin{tabular}{lccc}
    \textbf{For-comprehension}
      & \textbf{Rendezvous}
      & \textbf{Matching}
      & \textbf{Total} \\
    \midrule
    Orchestrator: $\mathtt{aliceDebit?!}$
      & 1 & 1 & 2 \\
    Orchestrator: $\mathtt{bobCredit?!}$
      & 1 & 1 & 2 \\
    Debit handler: $\forc$
      & 1 & 1 & 2 \\
    Credit handler: $\forc$
      & 1 & 1 & 2 \\
    \midrule
    \textbf{Total} & \textbf{4} & \textbf{4} & \textbf{8} \\
  \end{tabular}
\end{center}

\noindent
The six-token count from the sugar expansion is a
\emph{syntactic} count of explicitly written signed layers;
the eight-token count here is the \emph{semantic} count after
full desugaring of~$\mathtt{?!}$.  The static analysis of
Section~\ref{sec:linear-proof} must use the semantic count,
since it is the desugared form that the runtime executes.

The token stack must supply \emph{all} of these or
\emph{none} of them.  Partial supply leads to partial
execution, which is the vulnerability.

\subsection{Static Analysis as a Linear Resource Proof}
\label{sec:linear-proof}

The partial funding vulnerability is eliminated by a
\emph{static analysis} performed at deployment-acceptance
time --- the moment the validator's fuel gate fires and the
signed deployment enters execution.

\begin{definition}[Token demand]
  \label{def:token-demand}
  The \emph{token demand} of a signed term~$T$ with respect
  to a signature~$s$, written $\Delta_s(T)$, is the total
  number of $\signed{\cdots}{s}$ layers in~$T$ and in all
  processes reachable from~$T$ via synchronous calls ($?!$)
  or dequotation ($\drop{\cdot}$).  Formally:
  \begin{align*}
    \Delta_s(\signed{P}{s})
      &= 1 + \Delta_s(P) \\
    \Delta_s(\signed{P}{s'})
      &= \Delta_s(P)
      \qquad\text{for } s' \neq s \\
    \Delta_s(\recv{y}{x}{T})
      &= \Delta_s(T) \\
    \Delta_s(\send{x}{U})
      &= \Delta_s(U) \\
    \Delta_s(T \Par U)
      &= \Delta_s(T) + \Delta_s(U) \\
    \Delta_s(\drop{x})
      &= \Delta_s^{\mathit{resolve}}(x)
  \end{align*}
  where $\Delta_s^{\mathit{resolve}}(x)$ resolves the
  name~$x$ to its definition (if statically known) and
  computes $\Delta_s$ on the result.
\end{definition}

\begin{definition}[Token supply]
  \label{def:token-supply}
  The \emph{token supply} of a signature~$s$ in a system is
  the depth of $s$-indexed token stacks in the ambient
  parallel composition:
  \begin{align*}
    \Sigma_s(s : S)
      &= 1 + \Sigma_s(S) \\
    \Sigma_s(s' : S)
      &= \Sigma_s(S)
      \qquad\text{for } s' \neq s \\
    \Sigma_s(())
      &= 0 \\
    \Sigma_s(T \Par U)
      &= \Sigma_s(T) + \Sigma_s(U)
  \end{align*}
\end{definition}

\begin{definition}[Funding proof obligation]
  \label{def:funding-proof}
  A deployment is \emph{fully funded} if, for every
  signature~$s$ appearing in the deployment:
  \begin{equation}
    \Sigma_s \;\geq\; \Delta_s
    \label{eq:funding-obligation}
  \end{equation}
  This is a \emph{linear resource inequality}: each token is
  consumed exactly once, and the total supply must meet or
  exceed the total demand.
\end{definition}

\begin{theorem}[Decidability of the funding check]
  \label{thm:decidability}
  For a deployment whose call graph is statically known (no
  higher-order channel passing, no recursive dequotation
  chains), the funding proof
  obligation~\eqref{eq:funding-obligation} is decidable in
  time linear in the size of the deployment's AST.
\end{theorem}

\noindent
\textit{Proof sketch.}  $\Delta_s$ is computed by a single
pass over the AST, incrementing a counter at each
$\signed{\cdots}{s}$ node.  $\Sigma_s$ is computed by
inspecting the token stacks in the ambient parallel
composition.  The comparison is a single integer inequality.
In the presence of higher-order channels or recursive
dequotation, the analysis becomes a conservative
over-approximation: unresolvable $\drop{x}$ terms contribute
an ``unknown'' demand, and the validator rejects unless the
supply exceeds the known lower bound plus a configurable
safety margin.\qed

\subsection{The Validator's Acceptance Protocol}
\label{sec:acceptance-protocol}

The static analysis integrates into the validator's acceptance
protocol (Appendix~\ref{app:example}) as follows:

\begin{enumerate}[nosep]
  \item The validator receives a signed deployment
    $\signed{D}{c}$ and looks up the client's token stack.
  \item \textbf{Before} executing any part of the deployment,
    the validator computes $\Delta_c(D)$ by static analysis
    of the deployment's AST and call graph.
  \item The validator computes $\Sigma_c$ from the available
    token stack.
  \item If $\Sigma_c \geq \Delta_c$, the deployment is
    \emph{accepted}: the validator knows, by the linear
    resource proof, that every for-comprehension in the
    deployment's execution will be funded.  Every substitution
    that begins will complete.  No funds will be stranded.
  \item If $\Sigma_c < \Delta_c$, the deployment is
    \emph{rejected}: the validator returns an error
    (``insufficient phlogiston'') without executing any part
    of the deployment.  No state change occurs.  No tokens
    are consumed.
\end{enumerate}

\noindent
This protocol converts the database-level
atomicity-by-construction of individual rewrite rules into
\emph{financial-level} atomicity of multi-step transactions:
either the entire transaction is funded and executes to
completion, or it is rejected before any step fires.  The gap
between ``each step is atomic'' and ``the whole transaction is
atomic'' is bridged by the static analysis at acceptance time.

\subsection{Deployment Boundaries and Financial Atomicity}
\label{sec:deploy-boundaries}

The \emph{deployment} is the unit of financial-transaction
atomicity.  A deployment is a signed term submitted to a
validator for execution.  The validator's acceptance protocol
(Section~\ref{sec:acceptance-protocol}) either accepts the
entire deployment (the linear resource proof succeeds) or
rejects it entirely (the proof fails).  No partial acceptance
is possible.

This deployment-boundary semantics resolves several scenarios
that are problematic in conventional transaction processing:

\paragraph{Duplicate deployments.}
Suppose two copies of the same Alice~$\to$~Bob payment
deployment arrive at a validator, and Alice's token stack
contains exactly three tokens --- enough for one payment but
not two.  The validator processes each deployment as a
separate acceptance decision:

\begin{enumerate}[nosep]
  \item The first deployment arrives.  The validator computes
    $\Delta_{\mathit{Alice}} = 3$ and $\Sigma_{\mathit{Alice}}
    = 3$.  The linear proof succeeds: $3 \geq 3$.  The
    deployment is accepted.  Three tokens are committed.
  \item The second deployment arrives.  The validator computes
    $\Delta_{\mathit{Alice}} = 3$ and $\Sigma_{\mathit{Alice}}
    = 0$ (the tokens are already committed to the first
    deployment).  The linear proof fails: $0 < 3$.  The
    deployment is rejected.  No state change occurs.
\end{enumerate}

\noindent
The second deployment is rejected \emph{before any step
executes}.  There is no risk of a partial debit, no stranded
purse, no inconsistent state.  The linear resource proof
ensures that accepted deployments run to completion and
rejected deployments have no effect.

\paragraph{Simultaneous arrival.}
If the two deployments arrive at the same time, the validator
treats them as a \emph{single deployment} consisting of the
parallel composition of both terms:
\[
  \signed{D_1}{\mathit{Alice}} \Par \signed{D_2}{\mathit{Alice}}
\]
The static analysis computes
$\Delta_{\mathit{Alice}} = 3 + 3 = 6$ but
$\Sigma_{\mathit{Alice}} = 3$.  The linear proof fails:
$3 < 6$.  The combined deployment is rejected.  Neither
payment executes.

This is the correct behavior: there are not enough resources
to complete both financial transactions, so neither should
begin.  The linear resource proof detects this automatically,
without locks, priority queues, or rollback mechanisms.

\paragraph{Interleaved racing execution.}
In a na\"{\i}ve runtime that interleaves execution steps from
multiple deployments without checking resource bounds, the
following race is possible: deployment~1's debit fires
(consuming 2 tokens), then deployment~2's debit fires
(consuming the last token), and now neither deployment has
enough tokens for its credit step.  Both payments stall with
funds stranded in purses.

The linear resource proof eliminates this race entirely.
Acceptance is a \emph{gate}: a deployment passes through the
gate only if it carries a proof that all its resources are
available.  Once accepted, those resources are committed and
unavailable to other deployments.  There is no interleaving
of acceptance and execution --- acceptance is a single,
atomic decision that precedes all execution.

\begin{remark}[Deployment acceptance as linear proof search]
  The semantics is precise: a deployment is accepted
  \emph{if and only if} a linear proof of its resource usage
  exists.  The tokens consumed by the proof are
  \emph{linearly} committed --- they cannot be used by any
  other proof.  Two deployments competing for the same tokens
  are two proof obligations competing for the same linear
  hypotheses.  At most one can succeed.  This is not an
  implementation choice or a heuristic --- it is a theorem of
  linear logic.
\end{remark}

\begin{remark}[Connection to linear logic]
  The funding proof obligation~\eqref{eq:funding-obligation}
  is a judgment in the resource-sensitive fragment of linear
  logic.  Each token is a \emph{linear resource} that must
  be consumed exactly once.  The demand $\Delta_s$ is the
  number of linear hypotheses required; the supply $\Sigma_s$
  is the number of linear hypotheses available.  A successful
  funding check is a \emph{proof} in the linear resource
  logic that the deployment is well-resourced.

  This connection is not merely an analogy.  The OSLF functor
  (observer-generated logic) associated with the
  cost-accounted GSLT can be extended to generate a linear
  resource logic whose judgments \emph{are} funding proofs.
  The static analysis is then a proof search in this logic,
  and the validator is a proof checker.
\end{remark}

%% ══════════════════════════════════════════════════════════════════
\section{Conclusion}
\label{sec:conclusion}

We have presented a revision of the rho calculus in which cost
accounting is folded directly into the grammar and operational
semantics, eliminating the need for an external extension or a
translation layer.  The central design moves are: signatures as
names, cryptographic quoting as a second axis of reflection,
signed terms pervading the syntax, token stacks as first-class
terms, parametric polymorphism of parallel composition, and
genericity over the cryptographic backend.

The identification of the for-comprehension as the transaction
primitive --- with rendezvous and matching as the two funded
actions, the substitution as the witness, and application to the
continuation as the state update --- yields a precise theory of
atomicity at two levels.  \emph{Database-transaction atomicity}
(individual communications) is a theorem of the rewrite rules.
\emph{Financial-transaction atomicity} (multi-step operations
like payments) is guaranteed by a decidable static analysis at
deployment-acceptance time: a linear resource proof that
verifies every for-comprehension in the deployment's call
graph is funded before any step executes.  Deployment boundaries
provide the unit of financial atomicity, and the linear proof
serves as the concurrency control mechanism --- competing
deployments are competing proof obligations over the same
linear resources, and at most one can succeed.

The result is a calculus that embodies \emph{valuable friction}
at the syntactic level: every communicating process must engage
the cost layer, and the cost layer is rich enough to guarantee
transactional integrity.  Because the rho calculus is the formal
core of Rholang, this calculus serves directly as the design
specification for rearchitecting phlogiston accounting in
Rholang~1.4 and the F1R3FLY platform.

\bigskip
\begin{flushright}
  \textit{E Pluribus Potentia}
\end{flushright}

\appendix

%% ══════════════════════════════════════════════════════════════════
\section{Source-to-Source Translation to Pure Rho Calculus}
\label{app:translation}

This appendix defines a compositional source-to-source
translation from the cost-accounted rho calculus of
Section~\ref{sec:calculus} to the \emph{pure} rho calculus.
The translation is intended as a small-step development path:
one implements it as a compiler pass, and the existing Rholang
runtime handles the translated source without modification.

The translation is organized into five mutually recursive
functions:

\begin{center}
  \begin{tabular}{llll}
    \textbf{Function} & \textbf{Domain} & \textbf{Codomain}
      & \textbf{Purpose} \\
    \midrule
    $\mathcal{P}\lsem\cdot\rsem$
      & Processes   & Pure processes
      & Translate processes \\
    $\mathcal{N}\lsem\cdot\rsem$
      & Names       & Pure names
      & Translate names \\
    $\mathcal{T}\lsem\cdot\rsem$
      & Signed terms & Pure processes
      & Translate signed terms \\
    $\mathcal{K}\lsem\cdot\rsem$
      & Token stacks & Pure processes
      & Translate token stacks \\
    $\Sigma\lsem\cdot\rsem$
      & Signatures  & Pure names
      & Translate signatures \\
  \end{tabular}
\end{center}

\subsection{Signature Translation: $\Sigma\lsem\cdot\rsem$}
\label{app:sig-trans}

Signatures are mapped to names in the pure rho calculus.  Each
signature becomes a channel on which fuel tokens will be
communicated.

\begin{align}
  \Sigma\lsem g \rsem
    &= \quot{H_g}
    \quad\text{for a canonical process~$H_g$ encoding ground
      signature~$g \in G$}
    \label{eq:app-sig-ground}
  \\[4pt]
  \Sigma\lsem \hashf P \rsem
    &= \quot{H(\mathcal{P}\lsem P \rsem)}
    \quad\text{where~$H(\cdot)$ is a hash-encoding process
      constructor}
    \label{eq:app-sig-hash}
  \\[4pt]
  \Sigma\lsem s_1 \comp s_2 \rsem
    &= \quot{\bigl(
        \drop{\Sigma\lsem s_1 \rsem}
        \Par
        \drop{\Sigma\lsem s_2 \rsem}
      \bigr)}
    \label{eq:app-sig-comp}
\end{align}

\noindent
Equation~\eqref{eq:app-sig-comp} exploits reflection: the
name for a compound signature is the quote of the parallel
composition of the dequotations of the component names,
ensuring injectivity.

\begin{remark}[Ground-signature encoding]
  The canonical process~$H_g$ must be injective in~$g$: distinct
  ground signatures yield distinct processes.  In practice, $H_g$
  can be a process whose serialized form is the byte sequence of
  the public key or key hash~$g$.  Similarly, $H(\cdot)$ encodes
  the output of the cryptographic hash function applied to the
  serialized process tree.
\end{remark}

\subsection{Name Translation: $\mathcal{N}\lsem\cdot\rsem$}
\label{app:name-trans}

\begin{align}
  \mathcal{N}\lsem \quot{T} \rsem
    &= \quot{\mathcal{T}\lsem T \rsem}
    \label{eq:app-name-quote}
  \\[4pt]
  \mathcal{N}\lsem s \rsem
    &= \Sigma\lsem s \rsem
    \label{eq:app-name-sig}
\end{align}

\noindent
Quoted signed terms become quoted translations.  Signatures
become their signature-channel names.

\subsection{Token-Stack Translation: $\mathcal{K}\lsem\cdot\rsem$}
\label{app:token-trans}

A token stack becomes a chain of sends, one per layer, on the
corresponding signature channels.  Each send carries the
translation of the remaining stack as its payload.

\begin{align}
  \mathcal{K}\lsem () \rsem
    &= \nil
    \label{eq:app-tok-empty}
  \\[4pt]
  \mathcal{K}\lsem s : S \rsem
    &= \send{\Sigma\lsem s \rsem}{\mathcal{K}\lsem S \rsem}
    \label{eq:app-tok-cons}
\end{align}

\noindent
The token $s : S$ becomes a message on channel
$\Sigma\lsem s \rsem$ whose payload is the code of the
remaining balance~$\mathcal{K}\lsem S \rsem$.  This is the
\emph{fuel} that a signed process must consume before it can
act.

\subsection{Signed-Term Translation: $\mathcal{T}\lsem\cdot\rsem$}
\label{app:signed-trans}

Signed terms translate to pure processes.

\begin{align}
  \mathcal{T}\lsem T \Par U \rsem
    &= \mathcal{T}\lsem T \rsem \Par \mathcal{T}\lsem U \rsem
    \label{eq:app-st-par}
  \\[4pt]
  \mathcal{T}\lsem S \rsem
    &= \mathcal{K}\lsem S \rsem
    \label{eq:app-st-stack}
\end{align}

\noindent
For signed processes $\signed{P}{s}$, the translation wraps
the process in a \emph{fuel gate}: a for-comprehension on the
signature channel that must receive a token before the process
can proceed.  The consumed token's payload (the remaining
balance) is released into the continuation via dequotation.

\medskip
\noindent
\textbf{Atomic signature:}
\begin{equation}
  \mathcal{T}\lsem \signed{P}{s} \rsem
  \;=\;
  \recv{t}{\Sigma\lsem s \rsem}{
    \bigl(\drop{t} \Par \mathcal{P}\lsem P \rsem\bigr)
  }
  \label{eq:app-st-signed-atomic}
\end{equation}

\noindent
The process $\mathcal{P}\lsem P \rsem$ is blocked behind the
fuel gate~$\recv{t}{\Sigma\lsem s \rsem}{\cdots}$.  When a
token arrives, $t$ is bound to the code of the remaining
balance; $\drop{t}$ releases it; and the translated process
runs in parallel.

\medskip
\noindent
\textbf{Compound signature:}
\begin{equation}
  \mathcal{T}\lsem \signed{P}{s_1 \comp s_2} \rsem
  \;=\;
  \recv{t_1}{\Sigma\lsem s_1 \rsem}{
    \recv{t_2}{\Sigma\lsem s_2 \rsem}{
      \bigl(
        \drop{t_1} \Par \drop{t_2}
        \Par \mathcal{P}\lsem P \rsem
      \bigr)
    }
  }
  \label{eq:app-st-signed-compound}
\end{equation}

\noindent
Both component channels must supply fuel.  This generalizes
to $n$-ary composition $s_1 \comp s_2 \comp \cdots \comp s_n$
by nesting~$n$ fuel gates.

\subsection{Process Translation: $\mathcal{P}\lsem\cdot\rsem$}
\label{app:proc-trans}

\begin{align}
  \mathcal{P}\lsem \nil \rsem
    &= \nil
    \label{eq:app-p-nil}
  \\[4pt]
  \mathcal{P}\lsem \recv{y}{x}{T} \rsem
    &= \recv{y}{\mathcal{N}\lsem x \rsem}{
        \mathcal{T}\lsem T \rsem
      }
    \label{eq:app-p-recv}
  \\[4pt]
  \mathcal{P}\lsem \send{x}{U} \rsem
    &= \send{\mathcal{N}\lsem x \rsem}{
        \mathcal{T}\lsem U \rsem
      }
    \label{eq:app-p-send}
  \\[4pt]
  \mathcal{P}\lsem P \Par Q \rsem
    &= \mathcal{P}\lsem P \rsem
       \Par \mathcal{P}\lsem Q \rsem
    \label{eq:app-p-par}
  \\[4pt]
  \mathcal{P}\lsem \drop{x} \rsem
    &= \drop{\mathcal{N}\lsem x \rsem}
    \label{eq:app-p-drop}
\end{align}

\subsection{Infrastructure Processes}
\label{app:infrastructure}

When the granularity of token holding (combined vs.\ split)
does not match what the fuel gates expect, infrastructure
processes mediate the conversion.  These are deployed as
persistent system-level contracts.

\begin{definition}[Splitter]
  \label{def:app-splitter}
  Converts a combined token for~$s_1 \comp s_2$ into separate
  tokens for~$s_1$ and~$s_2$:
  \[
    \mathrm{Split}(s_1, s_2)
    \;=\;
    \recv{t}{\Sigma\lsem s_1 \comp s_2 \rsem}{
      \bigl(
        \send{\Sigma\lsem s_1 \rsem}{\nil}
        \Par
        \send{\Sigma\lsem s_2 \rsem}{\drop{t}}
      \bigr)
    }
  \]
\end{definition}

\begin{definition}[Joiner]
  \label{def:app-joiner}
  Converts separate tokens for~$s_1$ and~$s_2$ into a combined
  token for~$s_1 \comp s_2$:
  \[
    \mathrm{Join}(s_1, s_2)
    \;=\;
    \recv{t_1}{\Sigma\lsem s_1 \rsem}{
      \recv{t_2}{\Sigma\lsem s_2 \rsem}{
        \send{\Sigma\lsem s_1 \comp s_2 \rsem}{
          (\drop{t_1} \Par \drop{t_2})
        }
      }
    }
  \]
\end{definition}

\begin{remark}[Persistence]
  In production, splitters and joiners must be
  \emph{replicated} (persistent) processes.  The standard
  rho-calculus encoding of replication via self-referencing
  through reflection applies:
  \[
    !\,P
    \;\triangleq\;
    \recv{y}{\quot{D}}{(\drop{y} \Par P)}
    \Par
    \send{\quot{D}}{D}
    \qquad\text{where } D = \recv{y}{\quot{D}}{(\drop{y} \Par P)}
  \]
\end{remark}

\subsection{Verification Sketch}
\label{app:verification}

We verify that each rewrite rule of the cost-accounted
calculus is faithfully simulated by the translation.  Write
$\Vert\cdot\Vert$ for the system-level translation that
applies $\mathcal{T}\lsem\cdot\rsem$ to every signed-term
component and $\mathcal{K}\lsem\cdot\rsem$ to every token
stack, composing them in parallel.

\bigskip
\noindent
\textbf{Rule~1.}\;
$\signed{\recv{y}{x}{T} \Par \send{x}{U}}{s} \Par s : S
\;\red\; \subst{T}{\quot{U}}{y} \Par S$

\medskip\noindent
\textit{Translation of LHS:}
\begin{align*}
  &\recv{t}{\Sigma\lsem s \rsem}{
    \bigl(
      \drop{t}
      \Par \recv{y}{\mathcal{N}\lsem x \rsem}{
        \mathcal{T}\lsem T \rsem
      }
      \Par \send{\mathcal{N}\lsem x \rsem}{
        \mathcal{T}\lsem U \rsem
      }
    \bigr)
  }
  \\
  &\Par\;
  \send{\Sigma\lsem s \rsem}{\mathcal{K}\lsem S \rsem}
\end{align*}

\noindent
\textit{Step~1} --- fuel consumption on
$\Sigma\lsem s \rsem$, binding
$t \mapsto \quot{\mathcal{K}\lsem S \rsem}$:
\[
  \mathcal{K}\lsem S \rsem
  \Par
  \recv{y}{\mathcal{N}\lsem x \rsem}{
    \mathcal{T}\lsem T \rsem
  }
  \Par
  \send{\mathcal{N}\lsem x \rsem}{
    \mathcal{T}\lsem U \rsem
  }
\]

\noindent
\textit{Step~2} --- communication on
$\mathcal{N}\lsem x \rsem$, binding
$y \mapsto \quot{\mathcal{T}\lsem U \rsem}$:
\[
  \mathcal{K}\lsem S \rsem
  \Par
  \mathcal{T}\lsem T \rsem
    \bigl\{\quot{\mathcal{T}\lsem U \rsem}\big/y\bigr\}
\]

\noindent
By compositionality of the translation (the translation
commutes with substitution, Proposition~\ref{prop:subst-commute}),
this equals
$\mathcal{T}\lsem \subst{T}{\quot{U}}{y} \rsem \Par
\mathcal{K}\lsem S \rsem = \Vert \subst{T}{\quot{U}}{y} \Par S
\Vert$.\quad$\checkmark$

\bigskip
\noindent
\textbf{Rule~2.}\;
$\signed{\recv{y}{x}{T} \Par \send{x}{U}}{s_1 \comp s_2}
\Par s_1 : S \Par s_2 : S'
\;\red\; \subst{T}{\quot{U}}{y} \Par S \Par S'$

\medskip\noindent
The compound-signature fuel gate
\eqref{eq:app-st-signed-compound} nests two
$\mathbf{for}$-comprehensions on $\Sigma\lsem s_1 \rsem$ and
$\Sigma\lsem s_2 \rsem$.  The two token translations
$\mathcal{K}\lsem s_1 : S \rsem$ and
$\mathcal{K}\lsem s_2 : S' \rsem$ supply fuel to each.
After both fire and the data communication on
$\mathcal{N}\lsem x \rsem$ completes, the result is
$\mathcal{K}\lsem S \rsem \Par \mathcal{K}\lsem S' \rsem \Par
\mathcal{T}\lsem \subst{T}{\quot{U}}{y} \rsem
= \Vert \subst{T}{\quot{U}}{y} \Par S \Par S' \Vert$.
\quad$\checkmark$

\bigskip
\noindent
\textbf{Rule~3.}\;
$\signed{\recv{y}{x}{T} \Par \send{x}{U}}{s_1 \comp s_2}
\Par (s_1 \comp s_2) : S
\;\red\; \subst{T}{\quot{U}}{y} \Par S$

\medskip\noindent
The fuel gate listens on $\Sigma\lsem s_1 \rsem$ and
$\Sigma\lsem s_2 \rsem$, but the token lives on
$\Sigma\lsem s_1 \comp s_2 \rsem$.  The splitter
(Definition~\ref{def:app-splitter}) mediates:
\[
  \mathrm{Split}(s_1, s_2)
  \Par
  \send{\Sigma\lsem s_1 \comp s_2 \rsem}{
    \mathcal{K}\lsem S \rsem
  }
  \;\red\;
  \send{\Sigma\lsem s_1 \rsem}{\nil}
  \Par
  \send{\Sigma\lsem s_2 \rsem}{
    \mathcal{K}\lsem S \rsem
  }
\]
Now both fuel channels are populated and the reduction
proceeds as in Rule~2.  The residual $\nil$ from the
$s_1$-token is absorbed by structural equivalence.
\quad$\checkmark$

\bigskip
\noindent
\textbf{Rule~4.}\;
$\signed{\recv{y}{x}{T}}{s_1}
\Par \signed{\send{x}{U}}{s_2}
\Par s_1 : S \Par s_2 : S'
\;\red\; \subst{T}{\quot{U}}{y} \Par S \Par S'$

\medskip\noindent
Each signed process has an atomic fuel gate.  The
separately signed for-comprehension listens on
$\Sigma\lsem s_1 \rsem$; the separately signed send listens
on $\Sigma\lsem s_2 \rsem$.  Both tokens supply fuel
independently.  After fuel consumption:
\[
  \mathcal{K}\lsem S \rsem
  \Par
  \recv{y}{\mathcal{N}\lsem x \rsem}{
    \mathcal{T}\lsem T \rsem
  }
  \Par
  \send{\mathcal{N}\lsem x \rsem}{
    \mathcal{T}\lsem U \rsem
  }
  \Par
  \mathcal{K}\lsem S' \rsem
\]
The data communication on $\mathcal{N}\lsem x \rsem$
completes, yielding
$\Vert \subst{T}{\quot{U}}{y} \Par S \Par S' \Vert$.
\quad$\checkmark$

\bigskip
\noindent
\textbf{Rule~5.}\;
$\signed{\recv{y}{x}{T}}{s_1}
\Par \signed{\send{x}{U}}{s_2}
\Par (s_1 \comp s_2) : S
\;\red\; \subst{T}{\quot{U}}{y} \Par S$

\medskip\noindent
The fuel gates listen on $\Sigma\lsem s_1 \rsem$ and
$\Sigma\lsem s_2 \rsem$ separately, but the token is
combined on $\Sigma\lsem s_1 \comp s_2 \rsem$.  The splitter
mediates exactly as in Rule~3, producing separate tokens.
After fuel consumption and data communication, the result is
$\Vert \subst{T}{\quot{U}}{y} \Par S \Vert$.
\quad$\checkmark$

\subsection{Compositionality}
\label{app:compositionality}

The correctness of the verification rests on the following
property:

\begin{proposition}[Translation commutes with substitution]
  \label{prop:subst-commute}
  For all signed terms~$T$, $U$ and variables~$y$:
  \[
    \mathcal{T}\lsem T \rsem
    \bigl\{
      \quot{\mathcal{T}\lsem U \rsem} \big/ y
    \bigr\}
    \;=\;
    \mathcal{T}\lsem
      \subst{T}{\quot{U}}{y}
    \rsem
  \]
\end{proposition}

\noindent
The proof proceeds by structural induction on~$T$, with the
base case at~$\drop{x}$ using the semantic substitution
mechanism of the rho calculus.  The critical case is:
\[
  \drop{\mathcal{N}\lsem x \rsem}
  \bigl\{
    \quot{\mathcal{T}\lsem U \rsem} \big/ y
  \bigr\}
  \;=\;
  \begin{cases}
    \mathcal{T}\lsem U \rsem
      & \text{if } y \Neq \mathcal{N}\lsem x \rsem \\
    \drop{\mathcal{N}\lsem x \rsem}
      & \text{otherwise}
  \end{cases}
\]
which matches
$\mathcal{P}\lsem (\drop{x})\{\quot{U}/y\} \rsem$
by definition of semantic substitution.

\subsection{Implementation Notes}
\label{app:impl}

The translation is designed as a \textbf{compiler pass} in the
Rholang toolchain:

\begin{enumerate}[nosep]
  \item \textbf{Parse} the extended syntax (processes, signed
    terms, token stacks, parameterized signatures) into an
    extended AST.
  \item \textbf{Apply} $\mathcal{P}\lsem\cdot\rsem$ to the
    top-level process, which recursively invokes
    $\mathcal{T}\lsem\cdot\rsem$,
    $\mathcal{K}\lsem\cdot\rsem$,
    $\mathcal{N}\lsem\cdot\rsem$, and
    $\Sigma\lsem\cdot\rsem$.
  \item \textbf{Inject} persistent $\mathrm{Split}$ and
    $\mathrm{Join}$ infrastructure for every compound
    signature that appears in the source, unless the
    programmer has explicitly provided them.
  \item \textbf{Emit} pure rho calculus (Rholang~1.1) source.
    The existing runtime, reducer, and RSpace handle
    execution without modification.
\end{enumerate}

\noindent
This gives the platform immediate access to cost-accounted
semantics while the native implementation
(Section~\ref{sec:implications}) is developed in parallel.

%% ══════════════════════════════════════════════════════════════════
\section{Worked Example: Validator with Fee Extraction}
\label{app:example}

This appendix develops a complete worked example: a simplified
model of a validator node in the F1R3FLY network, expressed
first in the cost-accounted rho calculus and then translated
to pure rho via the compiler pass of Appendix~\ref{app:translation}.
The example exercises the bootstrap problem, token lookup,
error handling, and conversion of client tokens into validator
fees.

\subsection{Scenario and Conventions}
\label{app:ex-scenario}

Fix a ground signature scheme~$G$.  The actors are:

\begin{center}
  \begin{tabular}{lll}
    \textbf{Actor} & \textbf{Signature} & \textbf{Role} \\
    \midrule
    Validator & $v \in G$ & Accepts and executes deploys \\
    Client    & $c \in G$ & Submits a signed deploy \\
  \end{tabular}
\end{center}

\noindent
Since signatures are names, $v$ and $c$ double as channel
names.  We introduce five auxiliary channels (pure-rho names):

\begin{center}
  \begin{tabular}{ll}
    $\quot{DQ_v}$ & The validator's \emph{deploy queue} \\
    $\quot{F_v}$ & The validator's fee-collection channel \\
    $\quot{E}$   & The error-reporting channel \\
    $\quot{A_c}$ & The address where the client deposits tokens \\
    $\quot{W_v}$ & The validator's \emph{phlogiston wallet} \\
  \end{tabular}
\end{center}

\paragraph{Stake vs.\ phlogiston.}
It is important to distinguish the validator's \emph{stake}
from its \emph{phlogiston} (phlogistons).  The stake is
the quantity of $v$-tokens locked in the staking contract; it
determines the validator's weight in Casper CBC consensus and
is subject to slashing.  Phlogiston is a \emph{renewable
supply} drawn from the stake and deposited at the wallet
channel~$\quot{W_v}$.  The validator's handler consumes
phlogiston from~$\quot{W_v}$ to process client
deployments.

At each \emph{epoch boundary} (when the set of staked
validators is updated), the protocol replenishes the
validator's phlogiston wallet, provided the validator
has followed the rules of the protocol.  If the validator
commits a slashing offense, the entire stake and any
remaining phlogistons are moved to a private
unforgeable channel where they are held pending review and
adjudication.  After adjudication, the tokens may be
returned (if the offense was spurious), partially
redistributed (as a penalty), or burned.  Because the
unforgeable channel is private, the tokens are effectively
revoked: no process can access them without explicit
release by the adjudication contract.

In the example below, we model the phlogiston wallet
as a message on~$\quot{W_v}$ carrying a token stack.  The
validator handler draws tokens from this wallet rather than
directly from stake.

\subsection{Token Stack Decomposition}
\label{app:ex-decomp}

We adopt the following structural equivalence for token stacks,
which asserts that tokens are \emph{fungible and decomposable}:

\begin{equation}
  s : S
  \;\Seq\;
  (s : ()) \Par S
  \qquad\text{(token decomposition)}
  \label{eq:token-decomp}
\end{equation}

\noindent
Read: a multi-layer token stack is structurally equivalent to
the parallel composition of a single-token stack and the
remaining stack.  Repeated application decomposes
$s : s : s : ()$ into three parallel single-token stacks
$\bigl(s : ()\bigr) \Par \bigl(s : ()\bigr) \Par
\bigl(s : ()\bigr)$.

This equivalence is justified operationally: in the
source-to-source translation
(Appendix~\ref{app:translation}), the token stack
$s : s : s : ()$ translates to a \emph{nested} chain of
sends on $\Sigma\lsem s \rsem$, each carrying the remainder
as payload.  Decomposition corresponds to the standard
property that a communication can peel off one layer,
releasing the rest.

\subsection{Client-Side Assembly}
\label{app:ex-client}

The client prepares a deployment~$D$ (a process containing a
communication) and deposits three tokens at address~$\quot{A_c}$:

\begin{equation}
  \boxed{
    C
    \;\triangleq\;
    \send{\quot{A_c}}{c : c : c : ()}
    \;\Par\;
    \send{\quot{DQ_v}}{\signed{D}{c}}
  }
  \label{eq:client-assembly}
\end{equation}

\noindent
Here $\send{\quot{A_c}}{c : c : c : ()}$ deposits the
three-token stack at the client's token address, and
$\send{\quot{DQ_v}}{\signed{D}{c}}$ sends the signed deployment to
the validator's deploy queue~$\quot{DQ_v}$.

For concreteness, let~$D$ be a simple communication:
\begin{equation}
  D \;\triangleq\;
  \recv{z}{x}{T_{\mathit{app}}} \Par \send{x}{U_{\mathit{app}}}
  \label{eq:deploy-body}
\end{equation}
where $x$, $T_{\mathit{app}}$, and $U_{\mathit{app}}$ are the
client's application-level channel, continuation, and payload.

\subsection{Validator Handler}
\label{app:ex-validator}

The validator handler is a process that:
\begin{enumerate}[nosep]
  \item receives a signed deployment on the deploy
    queue~$\quot{DQ_v}$;
  \item looks up the client's tokens at address~$\quot{A_c}$;
  \item if tokens are present, extracts one token as a fee and
    composes the deployment with the remaining tokens;
  \item if no tokens are present, reports an error.
\end{enumerate}

\noindent
We define the handler using the uniform-signing sugar
(Definition~\ref{def:sugar-uniform}):

\begin{equation}
  \boxed{
    \mathit{VH}
    \;\triangleq\;
    \signed{
      \recv{\mathit{dep}}{\quot{DQ_v}}{
        \recv{\mathit{tok}}{\quot{A_c}}{
          \mathit{FeeExtract}
        }
      }
    }{v}
  }
  \label{eq:validator-handler}
\end{equation}

\noindent
The single outer $\signed{\cdot}{v}$ suffices because the
uniform-signing sugar~\eqref{eq:sugar-uniform} nests
recursively.  Expanding one level:
$\signed{\recv{\mathit{dep}}{\quot{DQ_v}}{\recv{\mathit{tok}}
{\quot{A_c}}{\mathit{FeeExtract}}}}{v}$ becomes
$\signed{\recv{\mathit{dep}}{\quot{DQ_v}}{\signed{\recv{\mathit{tok}}
{\quot{A_c}}{\mathit{FeeExtract}}}{v}}}{v}$, and the inner
$\signed{\recv{\mathit{tok}}{\quot{A_c}}
{\mathit{FeeExtract}}}{v}$ is itself sugar, expanding to
$\signed{\recv{\mathit{tok}}{\quot{A_c}}
{\signed{\mathit{FeeExtract}}{v}}}{v}$.  The fully desugared
form therefore has three $\signed{\cdot}{v}$ layers,
consuming three phlogistons.

The fee-extraction process $\mathit{FeeExtract}$ applies
token decomposition~\eqref{eq:token-decomp} to the received
token stack and splits it into a fee portion and a deployment
portion:

\begin{equation}
  \mathit{FeeExtract}
  \;\triangleq\;
  \send{\quot{F_v}}{c : ()}
  \;\Par\;
  \drop{\mathit{dep}} \;\Par\; (c : c : ())
  \label{eq:fee-extract}
\end{equation}

\noindent
Here:
\begin{itemize}[nosep]
  \item $\send{\quot{F_v}}{c : ()}$ sends one token (the fee)
    to the validator's fee-collection channel~$\quot{F_v}$.
  \item $\drop{\mathit{dep}}$ dequotes the received deployment,
    recovering the signed term~$\signed{D}{c}$.
  \item $(c : c : ())$ is the remaining two-token stack,
    placed in parallel with the deployment so that the
    cost-accounted rewrite rules can fire.
\end{itemize}

\noindent
The token decomposition~\eqref{eq:token-decomp} justifies the
split: when the validator receives the three-token stack via
$\recv{\mathit{tok}}{\quot{A_c}}{\cdots}$, the bound
name~$\mathit{tok}$ carries $\quot{c : c : c : ()}$.
Dequoting and decomposing:
\[
  c : c : c : ()
  \;\Seq\;
  (c : ()) \Par (c : c : ())
\]

\begin{remark}[Error handling]
  \label{rem:error}
  If the client fails to deposit tokens, the lookup
  $\recv{\mathit{tok}}{\quot{A_c}}{\cdots}$ simply
  \emph{blocks}: no message is available on~$\quot{A_c}$,
  and the fuel gate prevents execution.  This is the
  \emph{default} error mode --- the deployment is silently
  rejected by resource starvation.

  For proactive error reporting, the validator can race the
  token lookup against a timeout.  Using the standard
  select encoding in the rho calculus (a pair of
  for-comprehensions on the token channel and a timeout
  channel, with the first to fire pre-empting the other):
  \[
    \recv{\mathit{tok}}{\quot{A_c}}{
      \signed{\mathit{FeeExtract}}{v}
    }
    \;\Par\;
    \recv{\mathit{\_}}{\quot{\mathit{timeout}}}{
      \signed{
        \send{\quot{E}}{
          \signed{\mathit{ErrNoTokens}}{v}
        }
      }{v}
    }
  \]
  where a timeout process concurrently sends on
  $\quot{\mathit{timeout}}$ after a deadline.
\end{remark}

\subsection{Bootstrap: The Full System}
\label{app:ex-bootstrap}

The validator handler~$\mathit{VH}$ is itself a signed
process~$\signed{\cdots}{v}$.  By the cost-accounted rewrite
rules, it cannot execute unless $v$-tokens are present.
These tokens come from the validator's phlogiston
wallet~$\quot{W_v}$, which is replenished at each epoch
boundary from the validator's stake.

To connect the handler to the wallet, we wrap it in a
for-comprehension that draws phlogistons from
$\quot{W_v}$:

\begin{equation}
  \boxed{
    \mathit{VB}
    \;\triangleq\;
    \recv{\mathit{phlo}}{\quot{W_v}}{
      \bigl(
        \mathit{VH} \Par \drop{\mathit{phlo}}
      \bigr)
    }
  }
  \label{eq:validator-boot}
\end{equation}

\noindent
When a message carrying a token stack arrives on
$\quot{W_v}$, the variable~$\mathit{phlo}$ is bound to
the quoted stack; $\drop{\mathit{phlo}}$ dequotes it,
releasing the tokens into the ambient parallel composition
where they can be consumed by $\mathit{VH}$'s signed layers.

The full bootstrapped system is:

\begin{equation}
  \boxed{
    \Omega
    \;\triangleq\;
    \underbrace{
      \mathit{VB}
    }_{\text{validator bootstrap}}
    \;\Par\;
    \underbrace{
      \send{\quot{W_v}}{v : v : v : ()}
    }_{\text{phlogistons from wallet}}
    \;\Par\;
    \underbrace{
      C
    }_{\text{client assembly}}
  }
  \label{eq:full-system}
\end{equation}

\noindent
Expanding $\mathit{VB}$ and $C$ (with $\mathit{VH}$ in
sugared form per~\eqref{eq:validator-handler}):
\begin{align}
  \Omega
  \;=\;\;
  &\recv{\mathit{phlo}}{\quot{W_v}}{
    \bigl(
      \signed{
        \recv{\mathit{dep}}{\quot{DQ_v}}{
          \recv{\mathit{tok}}{\quot{A_c}}{
            \mathit{FeeExtract}
          }
        }
      }{v}
      \Par \drop{\mathit{phlo}}
    \bigr)
  }
  \notag\\
  &\Par\; \send{\quot{W_v}}{v : v : v : ()}
  \notag\\
  &\Par\; \send{\quot{A_c}}{c : c : c : ()}
  \;\Par\;
  \send{\quot{DQ_v}}{\signed{D}{c}}
  \label{eq:full-expanded}
\end{align}

\noindent
Recall that $\mathit{VH}$ uses the recursive uniform-signing
sugar (Definition~\ref{def:sugar-uniform}): the single outer
$\signed{\cdot}{v}$ expands recursively to three
$\signed{\cdot}{v}$ layers in the fully desugared form, one
per for-comprehension plus one for $\mathit{FeeExtract}$.

The validator's stake (locked in the staking contract) is
\emph{not} consumed directly.  The three phlogistons
$v : v : v : ()$ on~$\quot{W_v}$ are a renewable budget
drawn from stake, matched exactly to the three desugared
signed layers of the handler.

\subsection{Reduction in the Cost-Accounted Calculus}
\label{app:ex-reduction-ca}

We trace the reduction of~$\Omega$ using the cost-accounted
rewrite rules of Section~\ref{sec:rewrites}.  At each step
we identify which rule applies.

\bigskip
\noindent
\textbf{Step~0} --- The bootstrap process draws execution
tokens from the wallet (\textsc{Comm} on~$\quot{W_v}$):

The for-comprehension $\recv{\mathit{phlo}}{\quot{W_v}}{\cdots}$
communicates with $\send{\quot{W_v}}{v : v : v : ()}$, binding
$\mathit{phlo} \mapsto \quot{v : v : v : ()}$.  After
dequotation of $\drop{\mathit{phlo}}$, the system becomes:
\begin{align*}
  &\signed{
    \recv{\mathit{dep}}{\quot{DQ_v}}{
      \signed{
        \recv{\mathit{tok}}{\quot{A_c}}{
          \signed{\mathit{FeeExtract}}{v}
        }
      }{v}
    }
  }{v}
  \;\Par\; v : v : v : ()
  \\
  &\Par\; \send{\quot{A_c}}{c : c : c : ()}
  \;\Par\;
  \send{\quot{DQ_v}}{\signed{D}{c}}
\end{align*}

\noindent
The phlogistons are now in the ambient parallel
composition, available to the handler's signed layers.  Note
that this step is \emph{not} a cost-accounted rewrite: the
bootstrap for-comprehension on~$\quot{W_v}$ is an unsigned
process (a pure-rho communication), not a signed term.  It
is free infrastructure, analogous to the operating system's
process scheduler.

\bigskip
\noindent
\textbf{Step~1} --- The outermost signed handler consumes one
phlogiston and receives the signed deployment
(\textbf{Rule~1}):

The signed handler's outermost layer is
$\signed{\recv{\mathit{dep}}{\quot{DQ_v}}{\cdots}}{v}$, and the send
$\send{\quot{DQ_v}}{\signed{D}{c}}$ provides a message on~$\quot{DQ_v}$.  With
token~$v : v : v : ()$, decomposed as
$(v : ()) \Par (v : v : ())$:
\begin{align*}
  &\signed{
    \recv{\mathit{dep}}{\quot{DQ_v}}{\cdots} \Par \send{\quot{DQ_v}}{\signed{D}{c}}
  }{v}
  \;\Par\; v : ()
  \\
  &\quad\red\quad
  \cdots\bigl\{\quot{\signed{D}{c}}\big/\mathit{dep}\bigr\}
\end{align*}
After this step, $\mathit{dep}$ is bound to
$\quot{\signed{D}{c}}$ and one phlogiston is consumed.
The system becomes:
\begin{align*}
  &\signed{
    \recv{\mathit{tok}}{\quot{A_c}}{
      \signed{\mathit{FeeExtract}}{v}
    }
  }{v}
  \;\Par\; v : v : ()
  \\
  &\Par\; \send{\quot{A_c}}{c : c : c : ()}
\end{align*}

\noindent
(with $\mathit{dep} = \quot{\signed{D}{c}}$ captured in
$\mathit{FeeExtract}$).

\bigskip
\noindent
\textbf{Step~2} --- The second layer consumes one execution
token and looks up the client's tokens (\textbf{Rule~1}):
\begin{align*}
  &\signed{
    \recv{\mathit{tok}}{\quot{A_c}}{\cdots}
    \Par \send{\quot{A_c}}{c : c : c : ()}
  }{v}
  \;\Par\; v : ()
  \\
  &\quad\red\quad
  \signed{\mathit{FeeExtract}}{v}
    \bigl\{\quot{c : c : c : ()}\big/\mathit{tok}\bigr\}
  \;\Par\; v : ()
\end{align*}

\noindent
Now $\mathit{tok} = \quot{c : c : c : ()}$ and one more
phlogiston is consumed.  One phlogiston remains.

\bigskip
\noindent
\textbf{Step~3} --- The innermost layer consumes the last
phlogiston and executes fee extraction
(\textbf{Rule~1}):
\[
  \signed{\mathit{FeeExtract}}{v}
  \;\Par\; v : ()
  \;\;\red\;\;
  \mathit{FeeExtract} \Par ()
\]

\noindent
Substituting the bound values into~$\mathit{FeeExtract}$
\eqref{eq:fee-extract} and applying token decomposition:
\begin{equation}
  \send{\quot{F_v}}{c : ()}
  \;\Par\;
  \signed{D}{c}
  \;\Par\;
  (c : c : ())
  \label{eq:after-fee}
\end{equation}

\noindent
The system now contains the signed deployment~$\signed{D}{c}$
in parallel with two client tokens~$(c : c : ())$, plus a
fee token~$(c : ())$ en route to the validator's fee channel.

\bigskip
\noindent
\textbf{Step~4} --- The client deployment fires
(\textbf{Rule~1}):

Expanding $D = \recv{z}{x}{T_{\mathit{app}}} \Par
\send{x}{U_{\mathit{app}}}$:
\begin{align*}
  &\signed{
    \recv{z}{x}{T_{\mathit{app}}}
    \Par \send{x}{U_{\mathit{app}}}
  }{c}
  \;\Par\; c : c : ()
  \\[4pt]
  &\quad\red\quad
  \subst{T_{\mathit{app}}}
    {\quot{U_{\mathit{app}}}}{z}
  \;\Par\; c : ()
\end{align*}

\noindent
One client token is consumed.  The application continuation
runs with one remaining client token.

\bigskip
\noindent
\textbf{Final state:}
\[
  \send{\quot{F_v}}{c : ()}
  \;\Par\;
  \subst{T_{\mathit{app}}}
    {\quot{U_{\mathit{app}}}}{z}
  \;\Par\;
  (c : ())
\]

\noindent
The fee token is deposited at the validator's fee channel.
The application continuation runs.  One client token remains
for further computation.

\subsection{Source-to-Source Translation}
\label{app:ex-translation}

We now apply the translation of
Appendix~\ref{app:translation} to the system~$\Omega$.
Write $n_v = \Sigma\lsem v \rsem$ and
$n_c = \Sigma\lsem c \rsem$ for the translated signature
channels.

\medskip
\noindent
\textbf{Validator tokens:}
\begin{align*}
  \mathcal{K}\lsem v : v : v : () \rsem
  &= \send{n_v}{\send{n_v}{\send{n_v}{\nil}}}
\end{align*}

\noindent
A triple-nested send on~$n_v$.

\medskip
\noindent
\textbf{Client tokens:}
\begin{align*}
  \mathcal{K}\lsem c : c : c : () \rsem
  &= \send{n_c}{\send{n_c}{\send{n_c}{\nil}}}
\end{align*}

\medskip
\noindent
\textbf{Validator handler} (outermost fuel gate only,
for brevity):
\begin{align*}
  \mathcal{T}\lsem \mathit{VH} \rsem
  \;=\;\;
  &\recv{t_1}{n_v}{\bigl(
    \drop{t_1} \Par
    \recv{\mathit{dep}}{n_v}{\mathcal{T}\lsem\text{inner}\rsem}
  \bigr)}
\end{align*}

\noindent
where $\mathcal{T}\lsem\text{inner}\rsem$ nests two more
fuel gates on~$n_v$ (one per remaining $\signed{\cdots}{v}$
layer), followed by the token lookup on~$\quot{A_c}$, the
fee extraction, and the translated deployment.

\medskip
\noindent
\textbf{Fee extraction in the translated system:}

The critical observation is that fee extraction corresponds to
an \emph{additional fuel gate} on the client's signature
channel~$n_c$.  After the validator handler runs, the
translated system contains (among other terms):
\begin{equation}
  \underbrace{
    \recv{t_{\mathit{fee}}}{n_c}{
      \bigl(
        \send{\quot{F_v}}{\nil}
        \Par \drop{t_{\mathit{fee}}}
      \bigr)
    }
  }_{\text{fee interceptor}}
  \;\Par\;
  \underbrace{
    \send{n_c}{\send{n_c}{\send{n_c}{\nil}}}
  }_{\text{client's 3 tokens}}
  \label{eq:fee-interceptor}
\end{equation}

\noindent
The fee interceptor consumes the outermost send on~$n_c$
(one token).  Its payload is
$\send{n_c}{\send{n_c}{\nil}}$ --- the remaining two
tokens --- which is released via $\drop{t_{\mathit{fee}}}$.
The fee is recorded as a send on~$\quot{F_v}$.

After this reduction:
\[
  \send{\quot{F_v}}{\nil}
  \;\Par\;
  \send{n_c}{\send{n_c}{\nil}}
  \;\Par\;
  \underbrace{
    \recv{t}{n_c}{
      \bigl(\drop{t} \Par
        \mathcal{P}\lsem D \rsem
      \bigr)
    }
  }_{\text{deployment fuel gate}}
\]

\noindent
The deployment's own fuel gate now consumes from the
remaining two-token chain on~$n_c$, exactly as verified in
Appendix~\ref{app:verification}.

\subsection{Reduction Trace in Pure Rho}
\label{app:ex-trace-rho}

We trace the key reductions in the translated system,
suppressing intermediate structural equivalence steps.

\begin{enumerate}[label=\textbf{Step \arabic*.}, leftmargin=3em, start=0]
  \item \textbf{Wallet draw} on~$\quot{W_v}$.  The bootstrap
    for-comprehension receives the phlogiston stack.
    $\drop{\mathit{phlo}}$ releases three sends on~$n_v$
    into the ambient parallel composition.

  \item \textbf{Validator fuel gate~1} fires on~$n_v$.
    The outermost send in the phlogiston chain is
    consumed.  The handler's first $\forc$ unblocks.
    Remainder: two sends on~$n_v$.

  \item \textbf{Deployment received} on~$n_v$.  The
    translated $\send{n_v}{\mathcal{T}\lsem\signed{D}{c}\rsem}$
    communicates with the handler's inner $\forc$ on~$n_v$.
    $\mathit{dep}$ is bound.

  \item \textbf{Validator fuel gate~2} fires on~$n_v$.
    One more phlogiston consumed.

  \item \textbf{Token lookup} on~$\quot{A_c}$.  The
    client's token deposit communicates with the handler's
    $\forc$ on~$\quot{A_c}$.  $\mathit{tok}$ is bound
    to the quoted three-token chain.

  \item \textbf{Validator fuel gate~3} fires on~$n_v$.
    Last phlogiston consumed.  $\mathit{FeeExtract}$
    is now active.

  \item \textbf{Fee interceptor} fires on~$n_c$
    \eqref{eq:fee-interceptor}.  One client token consumed
    as fee.  Two-token chain released.

  \item \textbf{Deployment fuel gate} fires on~$n_c$.
    One client token consumed.  Application continuation
    $\mathcal{P}\lsem T_{\mathit{app}} \rsem$ unblocks.
    One-token chain released.

  \item \textbf{Application communication} on
    $\mathcal{N}\lsem x \rsem$.  The translated
    $\recv{z}{\mathcal{N}\lsem x \rsem}{\cdots}$ and
    $\send{\mathcal{N}\lsem x \rsem}{\cdots}$
    communicate.  The application payload is delivered.
\end{enumerate}

\noindent
\textbf{Final state} (in pure rho):
\[
  \underbrace{
    \send{\quot{F_v}}{\nil}
  }_{\text{fee deposited}}
  \;\Par\;
  \underbrace{
    \mathcal{T}\lsem
      \subst{T_{\mathit{app}}}
        {\quot{U_{\mathit{app}}}}{z}
    \rsem
  }_{\text{application result}}
  \;\Par\;
  \underbrace{
    \send{n_c}{\nil}
  }_{\text{1 client token remaining}}
\]

\subsection{Discussion}
\label{app:ex-discussion}

\paragraph{Stake, phlogistons, and epochs.}
The validator's \emph{stake} --- the quantity of $v$-tokens
locked in the staking contract --- is not consumed directly
by the handler.  Instead, the staking contract deposits a
renewable supply of \emph{phlogistons} at the wallet
channel~$\quot{W_v}$ at each epoch boundary (the point at
which the set of staked validators is updated).  The handler
draws from this wallet.

The phlogiston wallet must supply at least as many
tokens as the handler's nesting depth per deployment.  In the
example, three $\signed{\cdots}{v}$ layers require three
phlogistons per deploy.  A persistent validator
($!\,\mathit{VB}$) re-draws from the wallet after each
deployment, consuming phlogistons at a rate proportional
to the throughput of client deploys.

\paragraph{Slashing.}
In the current F1R3FLY slashing algorithm, any slashing
offense results in the removal of the validator's entire
stake.  In the cost-accounted model, slashing is implemented
as a \emph{channel revocation}: the slashing contract moves
the validator's stake and any remaining phlogistons to a
private unforgeable channel, where they are inaccessible to
the validator's handler.

Because the handler draws phlogistons from
$\quot{W_v}$, moving the wallet's contents to a private
channel immediately halts the validator: the bootstrap
for-comprehension $\recv{\mathit{phlo}}{\quot{W_v}}{\cdots}$
finds no message and blocks.  No further deployments can be
processed.  The validator is effectively offline.

The tokens held at the private channel are not destroyed.
They remain available for recovery pending review and
adjudication.  After adjudication, the tokens may be:
\begin{itemize}[nosep]
  \item \textbf{Returned} to the validator's wallet (if the
    offense was spurious or the result of a network fault).
  \item \textbf{Partially redistributed} to affected parties
    (as a penalty proportional to the severity of the
    offense).
  \item \textbf{Burned} (removed from circulation, reducing
    the total token supply).
\end{itemize}
The adjudication contract is itself a Rholang program,
expressible in the cost-accounted calculus, and subject to
the same formal verification tools.

\paragraph{Fee conversion.}
The fee token $(c : ())$ is a \emph{client-signature}
token.  For the validator to use it as fuel for its own
operations (which require $v$-signature tokens), the fee
must be \emph{converted}.  This conversion is a
market-making operation, expressible as a Rholang contract:
\[
  \mathrm{Exchange}(c, v)
  \;=\;
  \recv{t_c}{n_c}{
    \recv{t_v}{n_v}{
      \bigl(
        \send{n_c}{\drop{t_v}}
        \Par
        \send{n_v}{\drop{t_c}}
      \bigr)
    }
  }
\]
This persistent exchange contract swaps one $c$-token for
one $v$-token (and vice versa), implementing a 1:1 peg.
More sophisticated exchange contracts can implement
variable rates, order books, or automated market makers ---
all within Rholang.

\paragraph{Persistent validator.}
The example bootstrap processes a single deployment and
halts.  A production validator runs persistently,
re-drawing phlogistons from the wallet after each
deployment.  Using the standard reflective replication
encoding:
\[
  !\, \mathit{VB}
\]
Each iteration draws phlogistons from~$\quot{W_v}$,
processes a deployment, collects fees, and (via the exchange
contract) converts fees to replenish the phlogiston
supply --- establishing the economic feedback loop that
sustains the validator between epoch boundaries.

\paragraph{Connection to Casper CBC.}
In the F1R3FLY platform, the validator's stake
($v$-tokens locked in the staking contract) determines
the validator's weight in Casper CBC consensus.  The
cost-accounted calculus makes this relationship
\emph{intrinsic}: the validator literally cannot act
without phlogistons drawn from stake, and the rate of
token consumption is determined by the structure of the
handler --- the same structure that the consensus layer
inspects.  Slashing, epoch-boundary replenishment, and
stake-weighted voting are all expressible as
cost-accounted Rholang contracts, unifying the consensus
and execution layers under a single formal semantics.

%% ══════════════════════════════════════════════════════════════════
\section*{Acknowledgements}

The author thanks Mike Stay for extensive discussions on the
connections between linear logic and resource accounting,
Dylon Edwards for detailed feedback on the operational
semantics and implementation path, and Kevin Machiels for
insightful comments on the validator economics and slashing
model.

%% ══════════════════════════════════════════════════════════════════
\begin{thebibliography}{9}

\bibitem{MeredithRadestock2005}
  L.~G.~Meredith and M.~Radestock,
  ``A reflective higher-order calculus,''
  \textit{Electronic Notes in Theoretical Computer Science},
  vol.~141, no.~5, pp.~49--67, 2005.

\bibitem{Milner1999}
  R.~Milner,
  \textit{Communicating and Mobile Systems: the $\pi$-Calculus},
  Cambridge University Press, 1999.

\bibitem{RholangSpec}
  L.~G.~Meredith \textit{et al.},
  ``Rholang Specification,''
  F1R3FLY.io / RChain Cooperative, 2017--2026.

\bibitem{MeTTaIL}
  L.~G.~Meredith,
  ``MeTTaIL (Rholang 1.4): A language for defining languages,''
  F1R3FLY.io, 2025--2026.

\bibitem{VACRADA}
  L.~G.~Meredith \textit{et al.},
  ``Secure Blockchain-Based Health Data Platform: Cooperative
  Research and Development Agreement,''
  F1R3FLY Industries / U.S.\ Department of Veterans Affairs,
  2025--2026.

\end{thebibliography}

\end{document}
